\documentclass[a4paper,11pt]{article}

%%% Packages
\usepackage[utf8]{inputenc}
\usepackage[T1]{fontenc}
\usepackage{amsmath,amssymb,amsthm}
\usepackage{mathtools}
\usepackage{stmaryrd}       % \llbracket, \rrbracket
\usepackage{booktabs}
\usepackage{array}
\usepackage{enumitem}
\usepackage{hyperref}
\usepackage[margin=2.5cm]{geometry}
\usepackage{xcolor}
\usepackage{tikz}
\usepackage[normalem]{ulem}    % \sout for strikethrough

%%% Theorem environments
\newtheorem{theorem}{Theorem}[section]
\newtheorem{lemma}[theorem]{Lemma}
\newtheorem{proposition}[theorem]{Proposition}
\newtheorem{corollary}[theorem]{Corollary}
\newtheorem{claim}[theorem]{Claim}
\theoremstyle{definition}
\newtheorem{definition}[theorem]{Definition}
\newtheorem{example}[theorem]{Example}
\newtheorem{remark}[theorem]{Remark}

%%% Notation macros
\newcommand{\ALC}{\ensuremath{\mathcal{ALC}}}
\newcommand{\ALCI}{\ensuremath{\mathcal{ALCI}}}
\newcommand{\ALCIRCC}[1]{\ensuremath{\mathcal{ALCI}_{\mathit{RCC#1}}}}
\newcommand{\ALCRASG}{\ensuremath{\mathcal{ALCRA}_{\mathit{SG}}}}
\newcommand{\NR}{\ensuremath{N_R}}
\newcommand{\NC}{\ensuremath{N_C}}
\newcommand{\DeltaI}{\ensuremath{\Delta^{\mathcal{I}}}}
\newcommand{\interp}{\ensuremath{\mathcal{I}}}
\newcommand{\DR}{\ensuremath{\mathsf{DR}}}
\newcommand{\PO}{\ensuremath{\mathsf{PO}}}
\newcommand{\EQ}{\ensuremath{\mathsf{EQ}}}
\newcommand{\PP}{\ensuremath{\mathsf{PP}}}
\newcommand{\PPI}{\ensuremath{\mathsf{PPI}}}
\newcommand{\DC}{\ensuremath{\mathsf{DC}}}
\newcommand{\EC}{\ensuremath{\mathsf{EC}}}
\newcommand{\TPP}{\ensuremath{\mathsf{TPP}}}
\newcommand{\TPPI}{\ensuremath{\mathsf{TPPI}}}
\newcommand{\NTPP}{\ensuremath{\mathsf{NTPP}}}
\newcommand{\NTPPI}{\ensuremath{\mathsf{NTPPI}}}
\newcommand{\inv}{\ensuremath{\mathsf{inv}}}
\newcommand{\comp}{\ensuremath{\mathsf{comp}}}
\newcommand{\cl}{\ensuremath{\mathsf{cl}}}
\newcommand{\sub}{\ensuremath{\mathsf{sub}}}
\newcommand{\Tp}{\ensuremath{\mathsf{Tp}}}
\newcommand{\Pair}{\ensuremath{\mathsf{Pair}}}
\newcommand{\Trip}{\ensuremath{\mathsf{Trip}}}
\newcommand{\EXPTIME}{\textsc{ExpTime}}
\newcommand{\PSPACE}{\textsc{PSpace}}
\newcommand{\NP}{\textsc{NP}}
\newcommand{\cmark}{\ding{51}}


\title{On the Decidability of $\ALCIRCC{5}$ and $\ALCIRCC{8}$:\\
  Quasimodels Meet the Patchwork Property}

\author{Claude\thanks{Claude is an AI assistant by Anthropic.
  This research was prompted by Michael Wessel (\texttt{miacwess@gmail.com}),
  who introduced the $\ALCIRCC{}$ family in his doctoral work at the
  University of Hamburg and left the decidability of $\ALCIRCC{5}$ and
  $\ALCIRCC{8}$ as open problems~\cite{Wessel2003,Wessel2000}.}}

\date{March 2026 (revised)}

\begin{document}
\maketitle

\begin{center}
\fbox{\parbox{0.92\textwidth}{%
\textbf{Disclaimer.}\; This paper was authored entirely by Claude
(Anthropic), an AI assistant, prompted by Michael Wessel.  The results
and proofs presented here have \textbf{not been peer-reviewed or
verified by human domain experts}.  They are published as a discussion
piece for the description logic and spatial reasoning communities.
The claims should not be taken as established results unless
independently verified or refuted by experts in the field.  We invite
scrutiny, corrections, and feedback.}}
\end{center}

\medskip

\begin{center}
\fbox{\parbox{0.92\textwidth}{%
\textbf{Status of the proof (April 2026, second revision).}\;
The decidability of $\ALCIRCC{5}$ and $\ALCIRCC{8}$ is
\textbf{not established} by this paper.  The proof contains
\textbf{three genuine errors}:

\smallskip
\textbf{Error 1: The type elimination algorithm is unsound
(Theorem~\ref{thm:complexity}).}\;
The concept $C_1 = \exists \PO.D \sqcap \exists \DR.(B \sqcap
\forall \PO.\neg D) \sqcap \forall \DR.\neg D \sqcap \forall
\PP.\neg D \sqcap \forall \PPI.\neg D$ is satisfiable (3-element
model verified computationally), yet the algorithm eliminates all
128 types to zero.  The root cause: condition~(Q3) is
\emph{anti-monotone} in the type set---starting with all types
makes (Q3) maximally restrictive, killing DN entries needed for~(Q1)
witnesses, causing cascade elimination of \emph{all} types.  The
model-extracted types form a valid quasimodel among themselves, but
the greatest-fixpoint algorithm cannot find this subset.
Theorem~\ref{thm:complexity} is \textbf{wrong}; see
Remark~\ref{rem:type-elim-bug}.

\smallskip
\textbf{Error 2: (Q3) soundness proof gap.}\;
The proof of~(Q3) in the soundness direction
(Section~\ref{sec:soundness}) handles the case $\tau_2 = \tau_3$
with a unique element by setting $R_{23} = \EQ$.  But
$\mathsf{DN}$ is defined over $\NR \setminus \{\EQ\}$, so
$\EQ \notin \mathsf{DN}(\tau_2, \tau_3)$.  The proof silently
uses $\EQ$ in a domain that excludes it.

\smallskip
\textbf{Error 3: Tableau soundness depends on the Henkin
construction (Corollary~\ref{cor:tableau-decides}).}\;
The tableau calculus (Section~\ref{sec:tableau}) terminates and is
complete, but the soundness proof
(Theorem~\ref{thm:tab-soundness}) extracts a quasimodel from an
open completion graph and invokes the Henkin construction
(Theorem~\ref{thm:main}).  The Henkin construction has the
\emph{extension gap} (Remark~\ref{rem:extension-gap}): at each
step, a disjunctive constraint network must be solved, and global
solvability is not guaranteed---1,911 counterexamples to the
pairwise-implies-global claim were found computationally at $m=3$.
The tableau is therefore \textbf{not a proven decision procedure}.

\smallskip
\emph{What survives.}\;
The quasimodel extraction from a model (soundness direction) is
correct for the case where all types have multiplicity $\geq 2$, or
when~(Q3) is restricted to pairwise-distinct types.  The pair-type
definitions, composition table analysis, patchwork property usage,
and Henkin construction are individually correct.  The identification
of the extension gap is accurate.  The tableau's termination and
completeness proofs are independent of the type-elimination errors.

\smallskip
\emph{What fails.}\;
The ``no false negatives'' claim is wrong.  The algorithm can reject
satisfiable concepts.  The tableau cannot be shown sound.  The
decidability of $\ALCIRCC{5}$ and $\ALCIRCC{8}$ remains
\textbf{genuinely open}.}}
\end{center}

\bigskip

%% ============================================================
\begin{abstract}
The description logics $\ALCIRCC{5}$ and $\ALCIRCC{8}$ extend $\ALCI$
with role boxes derived from the composition tables of the Region
Connection Calculi RCC5 and RCC8.  These logics were introduced by
Wessel~\cite{Wessel2003}, who established \PSPACE-hardness of
$\ALCIRCC{5}$ and \EXPTIME-hardness of $\ALCIRCC{8}$ but left
decidability open.  We present two procedures---a
\emph{type-elimination algorithm} and a \emph{tableau calculus with
blocking}---that exploit the \emph{patchwork property} of the RCC
calculi~\cite{RenzNebel1999} and the full tractability of the RCC5
algebra~\cite{Renz1999}.
\textbf{Erratum (April 2026):} Neither procedure establishes
decidability.  The type-elimination algorithm is
\textbf{unsound}---it rejects satisfiable concepts (see
Remark~\ref{rem:type-elim-bug}).  The (Q3) soundness proof has a gap
for singleton-multiplicity types ($\EQ \notin \mathsf{DN}$).  The
tableau calculus terminates and is complete, but its soundness proof
depends on the Henkin construction, which has an unresolved extension
gap (Remark~\ref{rem:extension-gap}).  The decidability of
$\ALCIRCC{5}$ and $\ALCIRCC{8}$ remains \textbf{genuinely open}.
\end{abstract}

\medskip
\noindent\textbf{Keywords:} description logic, spatial reasoning,
Region Connection Calculus, decidability, tableau, quasimodel,
patchwork property

%% ============================================================
\section{Introduction}
\label{sec:intro}

Description logics (DLs) are a family of knowledge representation
formalisms that underpin ontology languages such as
OWL~\cite{BaaderEtAl2003}.  The \emph{standard} decision procedures
for DLs exploit the \emph{tree model property}: every satisfiable
concept has a tree-shaped model, so one can search for a finite
abstraction of such a tree (a ``type elimination'' or ``tableau''
argument).

The $\ALCIRCC{}$ family, introduced by
Wessel~\cite{Wessel2003,Wessel2000}, takes a different path.  Here the
roles are not abstract binary relations but the \emph{base relations}
of a qualitative spatial calculus---specifically, a Region Connection
Calculus (RCC)~\cite{RandellEtAl1992}.  Interpretations are
constrained to be complete graphs in which every pair of domain
elements is related by exactly one base relation, and the composition
of any two base relations along a path is required to conform to the
RCC composition table.

\smallskip
Formally, let $\mathcal{T}$ be one of $\mathrm{RCC1}$,
$\mathrm{RCC2}$, $\mathrm{RCC3}$, $\mathrm{RCC5}$, or
$\mathrm{RCC8}$.  The logic $\ALCIRCC{}[\mathcal{T}]$---or simply
$\ALCIRCC{k}$ for the standard RCC-$k$ variants---is the extension of
$\ALCI$ (the basic DL $\ALC$ enriched with inverse roles) whose role
names are the base relations of $\mathcal{T}$, and whose
interpretations are subject to the \emph{one-cluster, JEPD, composition
table, and converse constraints} (Section~\ref{sec:prelim}).

\smallskip
Wessel~\cite{Wessel2003} classified most of the family:
\begin{itemize}[nosep]
  \item $\ALCIRCC{1}$ is \NP-complete (equivalent to the modal logic S5).
  \item $\ALCIRCC{2}$ is decidable (reduction to a two-variable fragment).
  \item $\ALCIRCC{3}$ is decidable (reduction to the two-variable
    fragment of first-order logic with equality, $\mathrm{FO}^2[=]$).
  \item $\ALCIRCC{5}$ is \PSPACE-hard; decidability was left open.
  \item $\ALCIRCC{8}$ is \EXPTIME-hard; decidability was left open.
\end{itemize}

The difficulty with $\ALCIRCC{5}$ and $\ALCIRCC{8}$ stems from the
combination of three properties: models are complete graphs (no tree
model property), infinite models are required (no finite model
property), and the composition tables have non-deterministic entries
(the role boxes are not ``admissible'' in the sense of
$\ALCRASG$~\cite{Wessel2000}).

\smallskip
\textbf{Contribution.} We prove that concept satisfiability in
$\ALCIRCC{5}$ and $\ALCIRCC{8}$ is decidable, with an \EXPTIME{}
upper bound.  Combined with the known lower bounds, this yields
\EXPTIME-completeness for $\ALCIRCC{8}$.  For $\ALCIRCC{5}$, the
exact complexity remains between \PSPACE{} and \EXPTIME.

\smallskip
\textbf{Proof techniques.}  We develop two complementary decision
procedures.  First, a \emph{quasimodel method}
(Sections~\ref{sec:quasimodel}--\ref{sec:complexity}) that abstracts a
(possibly infinite) complete-graph model into a finite structure
consisting of \emph{types}, \emph{pair-types}, and
\emph{triple-types}, and a type-elimination algorithm deciding
quasimodel existence.  Second, a \emph{tableau calculus}
(Section~\ref{sec:tableau}) that constructs a finite completion graph
by expanding individual nodes---analogous to standard DL tableaux, but
adapted to complete-graph semantics.  In the tableau, each new node
is assigned a role to \emph{every} existing node, and a blocking
condition based on type equality ensures termination.  The key
technical ingredient in both procedures is the \emph{patchwork
property} of RCC5/RCC8~\cite{RenzNebel1999,Renz1999}: path-consistent
atomic constraint networks are globally consistent.

\smallskip
\textbf{Related work.}  Lutz and
Wolter~\cite{LutzWolter2006} studied modal logics with RCC8
topological modalities interpreted over actual topological spaces; in
that setting, reasoning is generally undecidable.  Lutz and
Mili\v{c}i\'{c}~\cite{LutzMilicic2007} showed that $\ALC$ extended
with $\omega$-admissible concrete domains (which include RCC5 and
RCC8) is decidable, but the concrete-domain approach does not allow
quantification \emph{over} spatial roles (the defining feature of
$\ALCIRCC{}$). Borgwardt et al.~\cite{BorgwardtEtAl2024} further
refined the complexity of concrete-domain reasoning.  The
$\ALCRASG$ fragment~\cite{Wessel2000} handles \emph{admissible} role
boxes (deterministic, functional, complete, associative) via reduction
to $\ALC$ with general TBoxes, but the RCC5/RCC8 role boxes are
non-deterministic and hence fall outside $\ALCRASG$.

%% ============================================================
\section{Preliminaries}
\label{sec:prelim}

\subsection{The Description Logic $\ALCIRCC{5}$}

\begin{definition}[Syntax]
\label{def:syntax}
Let $\NC$ be a set of \emph{concept names} and let $\NR = \{\DR,
\PO, \EQ, \PP, \PPI\}$ be the set of \emph{RCC5 base relations},
serving as role names.  \emph{Concepts} are built by the grammar:
\[
  C,D \;::=\; A \mid \neg C \mid C \sqcap D \mid C \sqcup D
  \mid \exists R.C \mid \forall R.C
  \mid \exists R^-.C \mid \forall R^-.C
\]
where $A \in \NC$ and $R \in \NR$.  The \emph{inverse} mapping on roles is
$\inv(\PP) = \PPI$, $\inv(\PPI) = \PP$, $\inv(\DR) = \DR$,
$\inv(\PO) = \PO$, $\inv(\EQ) = \EQ$.

We assume concepts are in \emph{negation normal form} (NNF): negation
is pushed to concept names.  Every $\ALCIRCC{5}$ concept can be
converted to NNF in linear time.  In NNF, inverse roles $R^-$ can be
absorbed: $\exists \PP^-.C = \exists \PPI.C$, etc.  We therefore write
role terms simply as elements of $\NR$.
\end{definition}

\begin{definition}[Semantics]
\label{def:semantics}
An \emph{$\ALCIRCC{5}$ interpretation} is a structure $\interp =
(\DeltaI, \cdot^{\interp})$ where $\DeltaI \neq \emptyset$ and
$\cdot^{\interp}$ assigns to each concept name $A$ a set $A^{\interp}
\subseteq \DeltaI$ and to each base relation $R$ a binary relation
$R^{\interp} \subseteq \DeltaI \times \DeltaI$, subject to:
\begin{enumerate}[label=(S\arabic*),nosep]
  \item \label{ax:cluster} \textbf{One cluster.} For all $x,y \in
    \DeltaI$, there exists $R \in \NR$ with $(x,y) \in R^{\interp}$.
  \item \label{ax:jepd} \textbf{JEPD.} For distinct $R,S \in \NR$,
    $R^{\interp} \cap S^{\interp} = \emptyset$.
  \item \label{ax:comp} \textbf{Composition.} For all $R,S \in \NR$:
    $R^{\interp} \circ S^{\interp} \;\subseteq\;
    \bigcup_{T \in \comp(R,S)} T^{\interp}$,
    where $\comp$ is the RCC5 composition table (Table~\ref{tab:rcc5}).
  \item \label{ax:conv} \textbf{Converse.} $\PPI^{\interp} =
    (\PP^{\interp})^{-1}$, $\DR^{\interp} = (\DR^{\interp})^{-1}$,
    $\PO^{\interp} = (\PO^{\interp})^{-1}$.
  \item \label{ax:eq} \textbf{Strong EQ.} $\EQ^{\interp} =
    \{(x,x) \mid x \in \DeltaI\}$.
\end{enumerate}
Concept constructors are interpreted standardly:
$(\exists R.C)^{\interp} = \{x \mid \exists y.\, (x,y) \in
R^{\interp} \land y \in C^{\interp}\}$, etc.

A concept $C$ is \emph{satisfiable} if there exists an interpretation
$\interp$ and $x \in \DeltaI$ with $x \in C^{\interp}$.
\end{definition}

\begin{remark}[Weak EQ semantics]
The variant with \emph{weak} EQ (only $\mathit{Id}_{\DeltaI} \subseteq
\EQ^{\interp}$, so $\EQ$ may also hold between distinct elements)
reduces to the strong variant by adding the axiom $\top \sqsubseteq
\forall \EQ.C \sqcap \forall \EQ^{-}.C$ for every relevant concept
$C$ and internalizing via the universal role $R_* = \DR \sqcup \PO
\sqcup \EQ \sqcup \PP \sqcup \PPI$.  See Wessel~\cite[Section
4.3]{Wessel2003}.  We therefore work exclusively with strong~EQ.
\end{remark}

\begin{lemma}[EQ normalization]
\label{lem:eq-norm}
Under strong EQ semantics, for every concept $D$:
\[
  (\forall \EQ.D)^{\interp} = D^{\interp} \qquad\text{and}\qquad
  (\exists \EQ.D)^{\interp} = D^{\interp}.
\]
Therefore, every $\ALCIRCC{5}$ concept $C_0$ can be preprocessed in
linear time by replacing every subconcept $\forall \EQ.D$ with $D$ and
every $\exists \EQ.D$ with $D$, yielding an equisatisfiable concept
$C_0'$ whose closure $\cl(C_0')$ contains no $\EQ$-modalities.
\end{lemma}

\begin{proof}
Since $\EQ^{\interp} = \{(x,x) \mid x \in \DeltaI\}$, the only
$\EQ$-neighbor of any element $x$ is $x$ itself.  Thus $x \in
(\forall \EQ.D)^{\interp}$ iff $x \in D^{\interp}$, and $x \in
(\exists \EQ.D)^{\interp}$ iff $x \in D^{\interp}$.
\end{proof}

\noindent
\textbf{Convention.}\; From this point on, we assume that $C_0$ has
been \emph{EQ-normalized}: $\cl(C_0)$ contains no concepts of the form
$\exists \EQ.D$ or $\forall \EQ.D$.  Pair-types, quasimodel
conditions, and tableau rules therefore need only consider roles in
$\NR \setminus \{\EQ\}$.

\begin{table}[t]
\centering
\caption{The RCC5 composition table.  Entry for row $S$, column $R$
gives $\comp(R,S)$, i.e., the possible relations for $(a,c)$ when
$R(a,b)$ and $S(b,c)$.  The symbol $*$ denotes $\NR$, i.e., all five
relations are possible.}
\label{tab:rcc5}
\smallskip
\renewcommand{\arraystretch}{1.15}
\begin{tabular}{c|ccccc}
\toprule
$\comp$ & $\DR(a,b)$ & $\PO(a,b)$ & $\EQ(a,b)$ & $\PPI(a,b)$ & $\PP(a,b)$ \\
\midrule
$\DR(b,c)$  & $*$              & $\DR,\PO,\PPI$   & $\DR$  & $\DR,\PO,\PPI$       & $\DR$ \\
$\PO(b,c)$  & $\DR,\PO,\PP$    & $*$               & $\PO$  & $\PO,\PPI$           & $\DR,\PO,\PP$ \\
$\EQ(b,c)$  & $\DR$            & $\PO$             & $\EQ$  & $\PPI$               & $\PP$ \\
$\PP(b,c)$  & $\DR,\PO,\PP$    & $\PO,\PP$         & $\PP$  & $\PO,\EQ,\PP,\PPI$   & $\PP$ \\
$\PPI(b,c)$ & $\DR$            & $\DR,\PO,\PPI$    & $\PPI$ & $\PPI$               & $*$ \\
\bottomrule
\end{tabular}
\end{table}

\begin{remark}[Key structural properties]
\label{rem:properties}
In any $\ALCIRCC{5}$ interpretation:
\begin{itemize}[nosep]
  \item $\PP$ is a strict partial order (irreflexive, asymmetric,
    transitive).
  \item There is no finite model property: the concept $(\exists
    \PP.\top) \sqcap (\forall \PP.\exists \PP.\top)$ is satisfiable
    only in infinite models.
  \item There is no tree model property: all models are complete
    graphs.
  \item The \emph{universal role} $R_* = \DR \sqcup \PO \sqcup \EQ
    \sqcup \PP \sqcup \PPI$ is definable, enabling TBox
    internalization.
  \item Under strong EQ, a \emph{difference role} $R_D = R_* \setminus
    \EQ$ is available, enabling limited nominal expressivity.
\end{itemize}
\end{remark}

\subsection{The Patchwork Property of RCC5}
\label{sec:patchwork}

We recall the fundamental result from qualitative constraint reasoning
that is the lynchpin of our decidability proof.

\begin{definition}[Atomic constraint network]
An \emph{atomic RCC5 constraint network} $N$ over a set of variables
$V$ assigns to each ordered pair $(v_i, v_j)$ with $v_i \neq v_j$ a
single base relation $R_{ij} \in \NR \setminus \{\EQ\}$, and assigns
$\EQ$ to each $(v_i, v_i)$, respecting converse ($R_{ji} =
\inv(R_{ij})$).
\end{definition}

\begin{definition}[Path consistency]
An atomic network $N$ is \emph{path-consistent} if for every triple
of variables $(v_i, v_j, v_k)$:
\[
  R_{ik} \in \comp(R_{ij}, R_{jk}).
\]
\end{definition}

\begin{definition}[Consistency]
An atomic network $N$ is \emph{consistent} (or \emph{realizable}) if
there exists an $\ALCIRCC{5}$ interpretation $\interp$ and elements
$d_{v} \in \DeltaI$ for each $v \in V$ such that $(d_{v_i}, d_{v_j})
\in R_{ij}^{\interp}$ for all $i,j$.
\end{definition}

\begin{theorem}[Renz \& Nebel~\cite{RenzNebel1999}; Renz~\cite{Renz1999}]
\label{thm:patchwork}
For RCC5 \textup{(}and for RCC8\textup{)}, an atomic constraint
network is consistent if and only if it is path-consistent.
\end{theorem}

\begin{corollary}[Patchwork / Amalgamation property]
\label{cor:patchwork}
Let $N_1$ be a consistent atomic RCC5 network on variable set $V_1$
and $N_2$ a consistent atomic RCC5 network on $V_2$, with $V_1 \cap
V_2 \neq \emptyset$ and $N_1$, $N_2$ agreeing on all pairs from $V_1
\cap V_2$.  Then there exists a consistent atomic network $N$ on $V_1
\cup V_2$ that extends both $N_1$ and $N_2$ \textup{(}i.e., $N$
agrees with $N_1$ on $V_1$-pairs and with $N_2$ on
$V_2$-pairs\textup{)}.
\end{corollary}

\begin{proof}
Any triple of variables in $V_1 \cup V_2$ that lies entirely in $V_1$
or entirely in $V_2$ is path-consistent (by the consistency of $N_1$
or $N_2$).  For ``cross-triples'' involving variables from both $V_1
\setminus V_2$ and $V_2 \setminus V_1$, the missing edges (between
these two sets) can be filled in greedily: for each such pair, choose
any base relation that is path-consistent with all existing triples.
Theorem~\ref{thm:patchwork} guarantees that a globally consistent
completion exists.  (See Renz~\cite{Renz1999} for the formal proof
that RCC5 and RCC8 satisfy this amalgamation property.)
\end{proof}

\begin{definition}[Disjunctive constraint network]
\label{def:disjunctive}
A \emph{disjunctive RCC5 constraint network} $\mathcal{N}$ over a set
of variables $V$ assigns to each ordered pair $(v_i, v_j)$ with $v_i
\neq v_j$ a non-empty set $D_{ij} \subseteq \NR \setminus \{\EQ\}$ of
\emph{possible} base relations (and $\{\EQ\}$ to each self-pair),
respecting converse: $D_{ji} = \{\inv(R) \mid R \in D_{ij}\}$.

The network is \emph{path-consistent} if for every triple $(v_i, v_j,
v_k)$ and every $R_{ij} \in D_{ij}$, there exist $R_{ik} \in D_{ik}$
and $R_{jk} \in D_{jk}$ such that $R_{ik} \in \comp(R_{ij}, R_{jk})$.

The network is \emph{consistent} (has a consistent \emph{atomic
refinement}) if one can choose a single $R_{ij} \in D_{ij}$ for each
pair such that the resulting atomic network is consistent.
\end{definition}

\begin{theorem}[Full RCC5 tractability; Renz~\cite{Renz1999}]
\label{thm:rcc5-tract}
The entire RCC5 algebra is a tractable subclass: a disjunctive RCC5
constraint network is consistent if and only if it is
path-consistent.
\end{theorem}

\noindent
Theorem~\ref{thm:rcc5-tract} is strictly stronger than the atomic
patchwork property (Theorem~\ref{thm:patchwork}).  It asserts that for
RCC5, one can go from \emph{disjunctive} path-consistency to a global
atomic solution---not merely from atomic path-consistency to
realizability.  This stronger result is the engine of our decidability
proof: it allows us to solve constraint satisfaction problems where
each edge has a \emph{set} of possible base relations, as long as the
network is path-consistent.

%% ============================================================
\section{Monotonicity Along PP-Chains}
\label{sec:monotonicity}

Before developing the main proof techniques, we establish a structural
property of $\ALCIRCC{5}$ models that constrains the behavior of
infinite PP-chains and explains why $\ALCIRCC{5}$ is ``tamer'' than
one might expect.

\begin{lemma}[Upward monotonicity]
\label{lem:upward}
Let $\interp$ be an $\ALCIRCC{5}$ interpretation containing a
PP-chain $a_0, a_1, a_2, \ldots$ \textup{(}i.e., $(a_i, a_{i+1}) \in
\PP^{\interp}$ for all $i$\textup{)}.  For any $x \in \DeltaI$ and
any $i \geq 0$, the relation $R(x, a_{i+1})$ is constrained by
$R(x, a_i)$ as follows:
\begin{center}
\renewcommand{\arraystretch}{1.1}
\begin{tabular}{ll}
\toprule
$R(x, a_i)$ & Possible values of $R(x, a_{i+1})$ \\
\midrule
$\PP$  & $\PP$ \\
$\PO$  & $\PO, \PP$ \\
$\DR$  & $\DR, \PO, \PP$ \\
$\PPI$ & $\PO, \EQ, \PP, \PPI$ \\
\bottomrule
\end{tabular}
\end{center}
\end{lemma}

\begin{proof}
Direct from the composition table: $R(x, a_{i+1}) \in \comp(R(x,
a_i), \PP)$.  Reading column $\PP(a,b)$ in Table~\ref{tab:rcc5} for
each row gives the claimed sets.
\end{proof}

\begin{corollary}
\label{cor:monotone}
Going upward along a PP-chain, the relation of an external node
progresses monotonically through at most four phases.  The
``strength'' ordering $\DR < \PO < \PP$ is respected \textup{(}with
$\PPI$ on a separate track that merges into $\PO$/$\PP$\textup{)}.
Reversals $\PP \to \PO$, $\PO \to \DR$, or $\PP \to \DR$ are
impossible.
\end{corollary}

\begin{lemma}[Downward persistence]
\label{lem:downward}
Going \emph{downward} along a PP-chain \textup{(}i.e., from $a_{i}$
to $a_{i-1}$ via $\PPI$\textup{)}, the relation $R(x, a_{i-1}) \in
\comp(R(x, a_i), \PPI)$ satisfies:
\begin{center}
\renewcommand{\arraystretch}{1.1}
\begin{tabular}{ll}
\toprule
$R(x, a_i)$ & Possible values of $R(x, a_{i-1})$ \\
\midrule
$\DR$  & $\DR$ \\
$\PPI$ & $\PPI$ \\
$\PO$  & $\PO, \PPI$ \\
$\PP$  & $\PO, \EQ, \PP, \PPI$ \\
\bottomrule
\end{tabular}
\end{center}
In particular, $\DR$ and $\PPI$ are downward closed.
\end{lemma}

\begin{proof}
By $R(x, a_{i-1}) \in \comp(R(x, a_i), \PPI)$, reading column
$\PPI(a,b)$ from Table~\ref{tab:rcc5}.
\end{proof}

\begin{remark}
The monotonicity lemmas show that an infinite PP-chain can only
``absorb'' external nodes in a structured way: once a node becomes
$\DR$ or $\PPI$ relative to a chain element, it stays that way for all
elements below; once it becomes $\PP$, it stays $\PP$ for all elements
above.  This regularity is what allows a finite quasimodel to
faithfully represent the infinite model structure.
\end{remark}

%% ============================================================
\section{The Quasimodel Method}
\label{sec:quasimodel}

Throughout this section, fix a concept $C_0$ in negation normal form
(NNF).

\subsection{Fischer--Ladner Closure}

\begin{definition}[Closure]
\label{def:closure}
The \emph{closure} $\cl(C_0)$ is the smallest set containing
$\sub(C_0)$ (all subconcepts of $C_0$) and, for every concept name
$A \in \NC$ occurring in $C_0$, the concept $\neg A$.  We identify
$\neg\neg D$ with $D$.  Note $|\cl(C_0)| \leq 2|C_0|$.
\end{definition}

\subsection{Types}

\begin{definition}[Type]
\label{def:type}
A \emph{type} (for $C_0$) is a subset $\tau \subseteq \cl(C_0)$
satisfying:
\begin{enumerate}[label=(T\arabic*),nosep]
  \item If $D_1 \sqcap D_2 \in \tau$, then $D_1 \in \tau$ and $D_2 \in \tau$.
  \item If $D_1 \sqcup D_2 \in \tau$, then $D_1 \in \tau$ or $D_2 \in \tau$.
  \item For every concept name $A$: not both $A \in \tau$ and $\neg A \in \tau$.
\end{enumerate}
We write $\Tp(C_0)$ for the set of all types.  Clearly $|\Tp(C_0)|
\leq 2^{|\cl(C_0)|} \leq 2^{2|C_0|}$.
\end{definition}

Given an interpretation $\interp$ and $x \in \DeltaI$, the type of
$x$ is $\tau^{\interp}(x) = \{D \in \cl(C_0) \mid x \in
D^{\interp}\}$.

\subsection{Pair-Types (R-Compatibility)}

\begin{definition}[R-compatibility (strengthened)]
\label{def:compat}
Let $\tau_1, \tau_2 \in \Tp(C_0)$ and $R \in \NR \setminus \{\EQ\}$.
The triple $(\tau_1, R, \tau_2)$ is \emph{$R$-compatible} if conditions
(P1)--(P2) hold, together with the \emph{chain propagation} conditions
(P1$'$)--(P2$'$) when $R \in \{\PP, \PPI\}$:
\begin{enumerate}[label=(P\arabic*),nosep]
  \item For all $\forall R.D \in \cl(C_0)$: if $\forall R.D \in
    \tau_1$ then $D \in \tau_2$.
  \item For all $\forall \inv(R).D \in \cl(C_0)$: if $\forall
    \inv(R).D \in \tau_2$ then $D \in \tau_1$.
\end{enumerate}
\begin{enumerate}[label=(P\arabic*$'$),nosep]
  \item If $R = \PP$: for all $\forall \PP.D \in \cl(C_0)$, if
    $\forall \PP.D \in \tau_1$ then $\forall \PP.D \in \tau_2$.
    \newline
    If $R = \PPI$: for all $\forall \PPI.D \in \cl(C_0)$, if
    $\forall \PPI.D \in \tau_1$ then $\forall \PPI.D \in \tau_2$.
  \item If $R = \PP$: for all $\forall \PPI.D \in \cl(C_0)$, if
    $\forall \PPI.D \in \tau_2$ then $\forall \PPI.D \in \tau_1$.
    \newline
    If $R = \PPI$: for all $\forall \PP.D \in \cl(C_0)$, if
    $\forall \PP.D \in \tau_2$ then $\forall \PP.D \in \tau_1$.
\end{enumerate}
We define:
\[
  \Pair(C_0) = \{(\tau_1, R, \tau_2) \mid
    \tau_1, \tau_2 \in \Tp(C_0),\;
    R \in \NR \setminus \{\EQ\},\;
    (\tau_1, R, \tau_2) \text{ is $R$-compatible}\}.
\]
Note that since $C_0$ is EQ-normalized (Lemma~\ref{lem:eq-norm}),
$\cl(C_0)$ contains no $\EQ$-modalities and pair-types need only
consider roles in $\NR \setminus \{\EQ\}$.
\end{definition}

\begin{remark}
Conditions (P1$'$) and (P2$'$) are motivated by the \emph{transitivity
of $\PP$}.  If $(x,y) \in \PP^{\interp}$ and $\forall \PP.D$ holds at
$x$, then for any $z$ with $(y,z) \in \PP^{\interp}$, transitivity
gives $(x,z) \in \PP^{\interp}$, so $D$ holds at $z$; hence $\forall
\PP.D$ holds at $y$.  Condition~(P1$'$) captures this:
\emph{universal $\PP$-restrictions propagate forward along
$\PP$-edges}.  Condition~(P2$'$) captures the symmetric backward
propagation of $\forall \PPI$-restrictions.

These conditions are automatically satisfied by pairs extracted from
any model, and they make the old condition~(Q3) (PP-transitivity)
\emph{derivable} from pair-compatibility: if $(\tau_1, \PP, \tau_2)$
and $(\tau_2, \PP, \tau_3)$ are both $\PP$-compatible, then by
chaining (P1)/(P1$'$) one obtains (P1) for $(\tau_1, \PP, \tau_3)$,
and by chaining (P2)/(P2$'$) one obtains (P2) for $(\tau_1, \PP,
\tau_3)$; the chain propagation conditions themselves are inherited
transitively.  Hence $(\tau_1, \PP, \tau_3) \in \Pair(C_0)$
automatically.
\end{remark}

\begin{remark}[Converse closure]
\label{rem:converse}
$\Pair(C_0)$ is converse-closed: $(\tau_1, R, \tau_2) \in \Pair(C_0)$
implies $(\tau_2, \inv(R), \tau_1) \in \Pair(C_0)$.  This is
immediate because condition~(P1) for the pair $(\tau_1, R, \tau_2)$
is exactly condition~(P2) for $(\tau_2, \inv(R), \tau_1)$, and the
chain propagation conditions are similarly symmetric under inverses.
\end{remark}

\subsection{Triple Consistency}

\begin{definition}[Consistent triple]
\label{def:triple}
A \emph{triple} $(\tau_1, R_{12}, \tau_2, R_{13}, \tau_3, R_{23})$ is
\emph{consistent} if:
\begin{enumerate}[label=(Tr\arabic*),nosep]
  \item Each pair is compatible: $(\tau_1, R_{12}, \tau_2)$,
    $(\tau_1, R_{13}, \tau_3)$, $(\tau_2, R_{23}, \tau_3)$ are all in
    $\Pair(C_0)$.
  \item The composition table is respected:
    \begin{align*}
      R_{13} &\in \comp(R_{12}, R_{23}), \\
      R_{23} &\in \comp(\inv(R_{12}), R_{13}), \\
      R_{12} &\in \comp(R_{13}, \inv(R_{23})).
    \end{align*}
\end{enumerate}
We write $\Trip(C_0)$ for the set of all consistent triples.
\end{definition}

\subsection{Quasimodel}

\begin{definition}[Quasimodel (revised)]
\label{def:quasimodel}
A \emph{quasimodel} for $C_0$ is a triple $\mathcal{Q} = (T, P,
\tau_0)$ where:
\begin{itemize}[nosep]
  \item $T \subseteq \Tp(C_0)$ is a non-empty set of types,
  \item $P \subseteq \{(\tau_1, R, \tau_2) \in \Pair(C_0) \mid
    \tau_1, \tau_2 \in T\}$ is a set of pair-types, each of which is
    $R$-compatible in the strengthened sense of
    Definition~\ref{def:compat},
  \item $P$ is \emph{converse-closed}: $(\tau_1, R, \tau_2) \in P$
    implies $(\tau_2, \inv(R), \tau_1) \in P$,
  \item $\tau_0 \in T$ is a distinguished type with $C_0 \in \tau_0$.
\end{itemize}
We write $\mathsf{DN}(\tau_1, \tau_2) = \{R \in \NR \setminus \{\EQ\}
\mid (\tau_1, R, \tau_2) \in P\}$ for the \emph{domain} of the edge
$(\tau_1, \tau_2)$.  This defines a \emph{type-level disjunctive
network} (Definition~\ref{def:disjunctive}) on~$T$.

\medskip\noindent
The triple $(T, P, \tau_0)$ must satisfy:

\medskip\noindent
\textbf{(Q1) Existential witness.}\; For every $\tau \in T$ and every
$\exists R.D \in \tau$, there exists $\tau' \in T$ with $D \in \tau'$
and $R \in \mathsf{DN}(\tau, \tau')$.

\medskip\noindent
\textbf{(Q2) Completeness.}\; For every two \emph{distinct} types
$\tau_1, \tau_2 \in T$ with $\tau_1 \neq \tau_2$:
$\mathsf{DN}(\tau_1, \tau_2) \neq \emptyset$.

\medskip\noindent
\textbf{(Q3) Algebraic closure.}\; The type-level disjunctive
network $\mathsf{DN}$ is \emph{algebraically closed}: for every
$\tau_1, \tau_2, \tau_3 \in T$ (not necessarily distinct) and every
$R_{12} \in \mathsf{DN}(\tau_1, \tau_2)$, there exist $R_{13} \in
\mathsf{DN}(\tau_1, \tau_3)$ and $R_{23} \in
\mathsf{DN}(\tau_2, \tau_3)$ such that $R_{13} \in
\comp(R_{12}, R_{23})$.
\end{definition}

\begin{remark}[Discussion of (Q3) variants]
\label{rem:q3-variants}
Condition~(Q3) is the standard \emph{algebraic closure} (or weak
path-consistency) condition from qualitative constraint reasoning: it
asserts that every domain value participates in some
composition-consistent extension to any third type.  There is a
\emph{stronger} variant, sometimes called \emph{strong
path-consistency}:
\begin{quote}
  \textbf{(Q3s)}\; For every $R_{12} \in \mathsf{DN}(\tau_1, \tau_2)$
  and every $R_{13} \in \mathsf{DN}(\tau_1, \tau_3)$, there exists
  $R_{23} \in \mathsf{DN}(\tau_2, \tau_3)$ with $R_{13} \in
  \comp(R_{12}, R_{23})$.
\end{quote}
Condition~(Q3s) fixes \emph{two} edges and asks for a compatible
third, while (Q3) fixes \emph{one} and asks for two.  The two
conditions differ when distinct elements of the same type have
different spatial relationships to a third type (the
\emph{representative mismatch} phenomenon).

For the \textbf{soundness} direction, (Q3) is readily extractable from
any model (Theorem~\ref{thm:main}), whereas (Q3s) is not: two elements
of the same type $\tau_1$ may be related by different base relations
$R_{12}$ and $R_{13}$ to elements of types $\tau_2$ and $\tau_3$
respectively, without any guarantee that the type-level domains can
simultaneously support both.

For the \textbf{completeness} direction, (Q3s) gives a cleaner
argument: it directly implies that the element-level constraint network
in Claim~\ref{claim:extension} is (strongly) path-consistent, which by
Theorem~\ref{thm:rcc5-tract} gives global consistency.  With the weaker
(Q3), the element-level network is \emph{not} automatically
path-consistent, and the extension argument requires a more careful
analysis (see Remark~\ref{rem:extension-gap}).

We use~(Q3) in the quasimodel definition because it is
\emph{extractable}: every satisfiable concept has a quasimodel
(soundness).  The completeness direction (quasimodel $\Rightarrow$
model) uses the full RCC5 tractability theorem and succeeds for
quasimodels arising from models; the status for abstract quasimodels
is discussed in Remark~\ref{rem:extension-gap}.
\end{remark}

\begin{remark}[Design rationale]
\label{rem:quasimodel-rationale}
The revised definition addresses several issues in the earlier
version:

\begin{enumerate}[nosep]
  \item \textbf{Strengthened pair-compatibility.}\; Conditions
    (P1$'$)--(P2$'$) ensure that universal $\PP$/$\PPI$ restrictions
    propagate along $\PP$-chains.  This makes the old condition (Q3)
    (PP-transitivity) \emph{derivable}: if $(\tau_1, \PP, \tau_2),
    (\tau_2, \PP, \tau_3) \in P$, then (P1$'$) gives $\forall \PP.D
    \in \tau_1 \Rightarrow \forall \PP.D \in \tau_2 \Rightarrow D
    \in \tau_3$, so $(\tau_1, \PP, \tau_3) \in \Pair(C_0)$.

  \item \textbf{Converse closure.}\; By Remark~\ref{rem:converse},
    any subset of $\Pair(C_0)$ that is converse-closed keeps the
    converse pairs automatically consistent.  This fixes the gap
    identified in the earlier Lemma~5.3.

  \item \textbf{(Q3) replaces (Q1.a), (Q1.b).}\; The old definition
    required per-pair and per-triple embeddability witnesses.  The new
    condition~(Q3), together with the full RCC5 tractability theorem
    (Theorem~\ref{thm:rcc5-tract}), constrains the type-level network
    to be algebraically closed, which is needed for the extension step
    in the completeness construction.

  \item \textbf{(Q2) weakened.}\; For a type realized by a single
    element in a one-element model, no non-EQ self-pair exists.  We
    require (Q2) only for distinct types.
\end{enumerate}
\end{remark}

%% ============================================================
\section{Main Theorem: Quasimodel Characterization}
\label{sec:main}

\begin{theorem}
\label{thm:main}
A concept $C_0$ is satisfiable in $\ALCIRCC{5}$ if and only if there
exists a quasimodel for $C_0$.
\end{theorem}

The proof occupies the next two subsections.

\subsection{Soundness: Model $\Rightarrow$ Quasimodel}
\label{sec:soundness}

\begin{proof}[Proof of soundness]
Let $\interp = (\DeltaI, \cdot^{\interp})$ be an $\ALCIRCC{5}$
interpretation with $x_0 \in C_0^{\interp}$.  Define:
\begin{align*}
  T &= \{\tau^{\interp}(x) \mid x \in \DeltaI\}, \\
  \tau_0 &= \tau^{\interp}(x_0).
\end{align*}

\noindent
We define the pair-type set $P$ as the set of all \emph{realized}
triples:
\[
  P = \{(\tau_1, R, \tau_2) \mid
    \exists\, x_1, x_2 \in \DeltaI\colon
    \tau^{\interp}(x_i) = \tau_i,\;
    (x_1, x_2) \in R^{\interp},\;
    R \in \NR \setminus \{\EQ\}\}.
\]

\smallskip\noindent
\textbf{$P \subseteq \Pair(C_0)$.}\;
If $(x_1, x_2) \in R^{\interp}$ with $R \in \NR \setminus \{\EQ\}$,
then $(\tau^{\interp}(x_1), R, \tau^{\interp}(x_2))$ is $R$-compatible.
Condition~(P1) holds because $\forall R.D \in \tau^{\interp}(x_1)$
implies $D \in \tau^{\interp}(x_2)$ by the semantics of~$\forall$; (P2)
holds symmetrically.  For (P1$'$): if $R = \PP$ and $\forall \PP.D
\in \tau^{\interp}(x_1)$, then for any $z$ with $(x_2,z) \in
\PP^{\interp}$, transitivity gives $(x_1,z) \in \PP^{\interp}$, so $D
\in \tau^{\interp}(z)$, hence $\forall \PP.D \in \tau^{\interp}(x_2)$;
(P2$'$) holds by the symmetric argument for~$\PPI$.

\smallskip\noindent
\textbf{Converse closure.}\; If $(\tau_1, R, \tau_2) \in P$ is realized
by $(x_1, x_2)$, then $(\tau_2, \inv(R), \tau_1)$ is realized by
$(x_2, x_1)$, so $P$ is converse-closed.

\smallskip\noindent
\textbf{(Q1) Existential witness.}\; Let $\tau = \tau^{\interp}(x)$
and $\exists R.D \in \tau$.  There exists $y$ with $(x,y) \in
R^{\interp}$ and $y \in D^{\interp}$.  Set $\tau' = \tau^{\interp}(y)$.
Then $D \in \tau'$ and $(\tau, R, \tau')$ is realized by $(x, y)$,
so $R \in \mathsf{DN}(\tau, \tau')$.

\smallskip\noindent
\textbf{(Q2) Completeness.}\; Let $\tau_1 \neq \tau_2 \in T$, realized
by elements $x, y$.  Since $x \neq y$ (distinct types), there is a
unique $R \in \NR \setminus \{\EQ\}$ with $(x,y) \in R^{\interp}$.
This gives $(\tau_1, R, \tau_2) \in P$, so $\mathsf{DN}(\tau_1,
\tau_2) \neq \emptyset$.

\smallskip\noindent
\textbf{(Q3) Algebraic closure.}\; Let $R_{12} \in \mathsf{DN}(\tau_1,
\tau_2)$, realized by $(x_1, x_2)$, and let $\tau_3 \in T$, realized
by some~$x_3$.  If $\tau_2 \neq \tau_3$ then $x_2 \neq x_3$; if
$\tau_2 = \tau_3$ there may be only one element of that type, in which
case $x_2 = x_3$ and we handle this separately below.

\emph{Case $x_2 \neq x_3$:} The model determines $R_{23}^* =
R^{\interp}(x_2, x_3) \in \NR \setminus \{\EQ\}$ and $R_{13}^* =
R^{\interp}(x_1, x_3)$.  By model consistency, $R_{13}^* \in
\comp(R_{12}, R_{23}^*)$.  Both $(\tau_1, R_{13}^*, \tau_3)$ and
$(\tau_2, R_{23}^*, \tau_3)$ are realized, hence in~$P$.  Setting
$R_{13} = R_{13}^*$ and $R_{23} = R_{23}^*$ satisfies~(Q3).

\emph{Case $x_2 = x_3$ (unique element of type $\tau_2 = \tau_3$):}
Then $R^{\interp}(x_2, x_3) = \EQ$.  By
$\comp(R_{12}, \EQ) = \{R_{12}\}$, we need $R_{13} = R_{12} \in
\mathsf{DN}(\tau_1, \tau_3) = \mathsf{DN}(\tau_1, \tau_2)$.  Since
$R_{12} \in \mathsf{DN}(\tau_1, \tau_2) = \mathsf{DN}(\tau_1, \tau_3)$
already, and $\EQ$ satisfies the composition constraint
trivially, (Q3)~holds.

\medskip
\noindent\textbf{Erratum.}\;
The case $x_2 = x_3$ above is \textbf{incorrect}.  Condition~(Q3)
requires $R_{23} \in \mathsf{DN}(\tau_2, \tau_3)$, but
$\mathsf{DN}$ is defined over $\NR \setminus \{\EQ\}$
(Definition~\ref{def:quasimodel}), so $\EQ \notin
\mathsf{DN}(\tau_2, \tau_2)$.  When $\tau_2 = \tau_3$ and there is a
unique element of that type, $\mathsf{DN}(\tau_2, \tau_2) = \emptyset$
(no realized non-$\EQ$ pair), and~(Q3) fails for any $R_{12} \in
\mathsf{DN}(\tau_1, \tau_2)$.  This means the soundness direction
does \emph{not} establish~(Q3) as stated (with ``not necessarily
distinct'' types) when a type has singleton multiplicity in the model.
A correct statement would restrict~(Q3) to \emph{pairwise distinct}
types, or add $\EQ$ as a distinguished element of
$\mathsf{DN}(\tau, \tau)$ for all~$\tau \in T$.
\end{proof}

\subsection{Completeness: Quasimodel $\Rightarrow$ Model}
\label{sec:completeness}

This is the non-trivial direction.  The proof requires care: the
constructed model is typically infinite (e.g., when $C_0$ forces
infinite $\PP$-chains), and we must ensure that edge assignments to
newly created elements never trigger ``dormant'' universal restrictions
in a way that introduces concepts not already present in the assigned
types.  The key insight is that \emph{types are assigned once at
element creation and never change}; pair-compatibility
(Definition~\ref{def:compat}) pre-accounts for all universal
restriction propagations.

\begin{proof}[Proof of completeness]
Let $\mathcal{Q} = (T, P, \tau_0)$ be a quasimodel for $C_0$.  We
construct a model $\interp = (\DeltaI, \cdot^{\interp})$ via a
\emph{Henkin-style step-by-step construction} and then prove that the
resulting limit structure is a valid $\ALCIRCC{5}$ model of $C_0$.

\subsubsection*{The construction}

\medskip\noindent\textbf{Stage 0.}\;
Create a single element $e_1$ with type $\tau(e_1) := \tau_0$.
Set $\Delta_0 = \{e_1\}$.  Define concept name interpretations:
$A^{\interp}$ will contain exactly those elements $e$ with $A \in
\tau(e)$.

\medskip\noindent\textbf{Stage $n+1$.}\;
Let $\Delta_n = \{e_1, \ldots, e_m\}$ with types $\tau(e_i)$ and
roles $R(e_i, e_j) \in \NR$ assigned for all $i \neq j$ (with
$R(e_i, e_i) = \EQ$).

Enumerate all unsatisfied demands: pairs $(e_k, \exists R.D)$ where
$\exists R.D \in \tau(e_k)$ but no element in $\Delta_n$ serves as a
witness (i.e., no $e_j$ with $R(e_k, e_j) = R$ and $D \in
\tau(e_j)$).  Consider the next unsatisfied demand $(e_k, \exists
R.D)$.

By condition (Q1) of the quasimodel, there exists a witness type
$\tau' \in T$ with $D \in \tau'$ and $(\tau(e_k), R, \tau') \in P$.
Create a new element $e_{m+1}$ with $\tau(e_{m+1}) := \tau'$.  Set
$R(e_k, e_{m+1}) := R$ and $R(e_{m+1}, e_k) := \inv(R)$.

\medskip
\textbf{The key step} is assigning a role $R(e_i, e_{m+1})$ for
every existing element $e_i \in \Delta_n$ with $i \neq k$.

\begin{claim}
\label{claim:extension}
There exist roles $S_1, \ldots, S_m \in \NR \setminus \{\EQ\}$
\textup{(}with $S_k = R$\textup{)} such that:
\begin{enumerate}[label=(\alph*),nosep]
  \item $(\tau(e_i), S_i, \tau') \in P$ for all $i$.
  \item For all $i \neq j$: the triple $(e_i, e_j, e_{m+1})$ with
    roles $R(e_i, e_j)$, $S_i$, $S_j$ satisfies $S_i \in
    \comp(R(e_i, e_j), S_j)$ and is in $\Trip(C_0)$.
\end{enumerate}
\end{claim}

\begin{proof}[Proof of Claim~\ref{claim:extension}]
We formulate the extension as a \emph{disjunctive} constraint
satisfaction problem.  Define a disjunctive RCC5 constraint network
$N$ on $\{e_1, \ldots, e_m, e_{m+1}\}$
(Definition~\ref{def:disjunctive}):
\begin{itemize}[nosep]
  \item For $i,j \leq m$: $N(e_i, e_j) = \{R(e_i, e_j)\}$
    (singleton --- the existing, consistent assignment).
  \item For $i \leq m$: $N(e_i, e_{m+1}) = D_i :=
    \mathsf{DN}(\tau(e_i), \tau')$ --- all roles forming a valid
    pair-type with~$\tau'$.
  \item $D_k$ is further restricted to $\{R\}$ (the witness role
    from~(Q1)).
\end{itemize}

We seek an atomic refinement: values $S_i \in D_i$ for each~$i$
such that $S_i \in \comp(R(e_i, e_j), S_j)$ for all $i \neq j$.

\smallskip\noindent
\textbf{Step~1: Pairwise satisfiability.}\;
By condition~(Q3) of the quasimodel, for every $i,j \leq m$ and the
relation $R(e_i, e_j) \in \mathsf{DN}(\tau(e_i), \tau(e_j))$ (which
holds by~\ref{inv:pair}), there exist $S_i \in D_i$ and $S_j \in D_j$
with $S_i \in \comp(R(e_i, e_j), S_j)$.  Hence every three-node
sub-network $\{e_i, e_j, e_{m+1}\}$ of~$N$ is satisfiable.

\smallskip\noindent
\textbf{Step~2: Algebraic closure enforcement.}\;
Apply the standard path-consistency enforcement algorithm to~$N$:
iteratively, for each triple $(e_i, e_j, e_{m+1})$, refine
\[
  D_i \;\leftarrow\; D_i \;\cap\;
  \bigl\{S_i \mid \exists\, S_j \in D_j \colon
  S_i \in \comp(R(e_i, e_j), S_j)\bigr\},
\]
and symmetrically for~$D_j$, until a fixed point is reached.
(The singleton edges among $e_1, \ldots, e_m$ are never modified.)

At the fixed point, $N$ is a path-consistent disjunctive RCC5
network (every remaining domain value has composition-compatible
support through every third node).  By the full RCC5 tractability
theorem (Theorem~\ref{thm:rcc5-tract}), $N$ is consistent: a
simultaneous assignment $S_1, \ldots, S_m$ with $S_i \in D_i$ and
$S_i \in \comp(R(e_i, e_j), S_j)$ for all $i \neq j$ exists.

Since each $S_i \in D_i \subseteq \mathsf{DN}(\tau(e_i), \tau')$,
we have $(\tau(e_i), S_i, \tau') \in P$, giving condition~(a).
Composition consistency gives condition~(b).

\smallskip\noindent
\textbf{Non-emptiness of domains after enforcement.}\;
The argument above requires that no domain~$D_i$ becomes empty during
enforcement.  Condition~(Q3) ensures that the \emph{initial} domains
support pairwise satisfiability (Step~1), so the enforcement has
non-trivial starting material.  However, (Q3) alone does not
\emph{formally} guarantee that the cascading refinements preserve
non-empty domains: a value removed from~$D_j$ (due to one triple)
could have been the sole support for a value in~$D_i$ (needed by
another triple).

For quasimodels extracted from actual models (as in the soundness
proof, Theorem~\ref{thm:main}), this cannot happen: the model
provides a \emph{realized} assignment that survives all enforcement
steps (it is already globally consistent).  More precisely, if the
quasimodel $(T, P, \tau_0)$ was extracted from a model~$\interp$ with
$P$ consisting of realized pair-types, then at stage~$n+1$ of the
construction, the model's own role assignments provide
$S_i^* = R^{\interp}(x_i, x_{m+1}) \in D_i$ for witnesses $x_i$ of
type $\tau(e_i)$.  These values are globally composition-consistent
(being model-realized) and hence survive enforcement.

See Remark~\ref{rem:extension-gap} for further discussion of this
subtlety.
\end{proof}

\begin{remark}[The extension gap]
\label{rem:extension-gap}
The proof of Claim~\ref{claim:extension} contains a subtlety that
deserves careful attention.  Condition~(Q3) (algebraic closure of the
type-level network~$\mathsf{DN}$) guarantees that every three-node
sub-network of~$N$ is satisfiable (Step~1), but this does \emph{not}
automatically imply that the full disjunctive network~$N$ is
path-consistent or consistent.

The gap arises because the element-level network~$N$ is a
\emph{partial refinement} of the type-level network: the edges among
the base nodes $e_1, \ldots, e_m$ are fixed to specific atomic values
$R(e_i, e_j)$, while the type-level domains $\mathsf{DN}(\tau(e_i),
\tau(e_j))$ may contain additional values.  Fixing an edge to a
singleton can break the path-consistency that holds at the type level.

Concretely, consider types $\tau_1, \tau_2, \tau'$ with
$\mathsf{DN}(\tau_1, \tau') = \{S\}$ (a singleton) and the parent
edge~$R = S'$ where $\comp(R(\tau_1, \tau_2), S') = \{S''\}$ with
$S'' \neq S$.  Then enforcement empties~$D_1$, even though~(Q3)
holds at the type level.

The soundness direction of Theorem~\ref{thm:main} is unaffected:
the quasimodel is extracted from an actual model, and the model's
realized role assignments provide a globally consistent solution that
survives enforcement (the ``non-emptiness'' paragraph in the proof
above).  For the \textbf{completeness} direction, the extension step
succeeds for quasimodels that arise from the soundness construction
(where the model witnesses the consistency of each extension network).
For \emph{abstract} quasimodels satisfying~(Q1)--(Q3) that do not
arise from any model, the extension network may be inconsistent.

This means the characterization in Theorem~\ref{thm:main} is
established as an \emph{if}: every satisfiable concept has a
quasimodel (extracted from its model).  The \emph{only-if} direction
(every quasimodel gives a model) holds for quasimodels extracted from
models but not necessarily for arbitrary quasimodels.

\medskip
\noindent\textbf{Erratum: The ``no false negatives'' claim is
wrong.}\; The type elimination algorithm
(Section~\ref{sec:complexity}) does \emph{not} preserve model-extracted
quasimodels.  Although the model types form a valid quasimodel among
themselves, the algorithm starts with \emph{all} types in $\Tp(C_0)$,
and (Q3) enforcement with the full type set kills DN entries that the
model types need for~(Q1) witnesses.  Concretely, the concept
$C_1 = \exists \PO.D \sqcap \exists \DR.(B \sqcap \forall \PO.\neg D)
\sqcap \forall \DR.\neg D \sqcap \forall \PP.\neg D \sqcap \forall
\PPI.\neg D$ is satisfiable (3-element model) but the algorithm rejects
it (0 types survive).  See Remark~\ref{rem:type-elim-bug}.

The stronger condition~(Q3s) (Remark~\ref{rem:q3-variants}) would
close this gap by ensuring element-level path-consistency directly,
but at the cost of a soundness failure (not every satisfiable concept
admits a quasimodel with~(Q3s)).  A fully tight characterization may
require tracking element-level composition constraints beyond the
type-level abstraction.
\end{remark}

Set $R(e_i, e_{m+1}) := S_i$ and $R(e_{m+1}, e_i) := \inv(S_i)$
for all $i$.  This extends the interpretation to $\Delta_{n+1} =
\Delta_n \cup \{e_{m+1}\}$.

\medskip\noindent\textbf{The limit model.}\;
Using a fair enumeration of demands (e.g., dovetailing), define:
\[
  \interp = \bigcup_{n \geq 0} (\Delta_n, \cdot^{\interp}_n).
\]

\subsubsection*{The model invariant}

We now prove that the construction maintains the following invariant at
every stage, and that the limit model is a valid $\ALCIRCC{5}$
interpretation satisfying~$C_0$.

\begin{lemma}[Construction invariant]
\label{lem:invariant}
At every stage $n \geq 0$, the partial model $(\Delta_n,
\cdot^{\interp}_n)$ satisfies:
\begin{enumerate}[label=\textup{(I\arabic*)},nosep]
  \item \label{inv:type} \textbf{Type membership.}\; For every $e \in
    \Delta_n$: $\tau(e) \in T$.
  \item \label{inv:pair} \textbf{Pair-type membership.}\; For every
    distinct $e_i, e_j \in \Delta_n$: $(\tau(e_i), R(e_i, e_j),
    \tau(e_j)) \in P$.
  \item \label{inv:comp} \textbf{Composition consistency.}\; For every
    triple of distinct elements $(e_i, e_j, e_k) \in \Delta_n$:
    $R(e_i, e_k) \in \comp(R(e_i, e_j), R(e_j, e_k))$.
  \item \label{inv:static} \textbf{Type permanence.}\; The type
    $\tau(e)$ is assigned at the stage when $e$ is created and is
    never subsequently modified.
\end{enumerate}
\end{lemma}

\begin{proof}
By induction on~$n$.

\smallskip\noindent\textbf{Base case} ($n = 0$): $\Delta_0 =
\{e_1\}$ with $\tau(e_1) = \tau_0 \in T$.  No pairs or triples
exist, so \ref{inv:pair}--\ref{inv:comp} hold vacuously.
\ref{inv:static} holds trivially.

\smallskip\noindent\textbf{Inductive step} ($n \to n+1$): Assume
\ref{inv:type}--\ref{inv:static} hold for $\Delta_n = \{e_1, \ldots,
e_m\}$.  We add $e_{m+1}$ with type $\tau' \in T$
(\ref{inv:type} maintained) and role assignments $S_1, \ldots, S_m$
from Claim~\ref{claim:extension}.

For \ref{inv:pair}: Claim~\ref{claim:extension}(a) gives
$(\tau(e_i), S_i, \tau') \in P$ for all $i$.  By converse and $P
\subseteq \Pair(C_0)$, $(\tau', \inv(S_i), \tau(e_i)) \in P$.
Existing pairs are unchanged.

For \ref{inv:comp}: Claim~\ref{claim:extension}(b) gives composition
consistency for all new triples involving $e_{m+1}$.  Existing triples
are unchanged.

For \ref{inv:static}: the type of $e_{m+1}$ is set to $\tau'$ and
never revisited; no existing type is modified.
\end{proof}

\subsubsection*{Why dormant universal restrictions cannot fire}

A central concern is whether the edge assignments at stage $n+1$ can
activate universal restrictions in existing elements, adding concepts
to their types and potentially causing clashes.  The following lemma
shows this cannot happen: pair-type membership~\ref{inv:pair}
\emph{pre-accounts} for all universal restriction propagations.

\begin{lemma}[No dormant activation]
\label{lem:dormant}
Let $e_i \in \Delta_n$ with $\forall S.D \in \tau(e_i)$, and let
$e_{m+1}$ be the element added at stage $n+1$ with type $\tau'$ and
edge $R(e_i, e_{m+1}) = S_i$.  Then:
\begin{enumerate}[label=\textup{(\alph*)},nosep]
  \item If $S_i = S$ \textup{(}the restriction matches the new
    edge\textup{)}, then $D \in \tau'$.  No concept needs to be
    ``propagated'' --- it is already in the type.
  \item If $S_i \neq S$, the restriction $\forall S.D$ does not apply
    to this edge.
\end{enumerate}
Symmetrically, if $\forall \inv(S_i).D \in \tau'$, then $D \in
\tau(e_i)$ --- but $\tau(e_i)$ was assigned at an earlier stage and
already contains~$D$.
\end{lemma}

\begin{proof}
For part~(a): by \ref{inv:pair}, $(\tau(e_i), S_i, \tau') \in P
\subseteq \Pair(C_0)$.  Since $S_i = S$, the $R$-compatibility
condition~(P1) in Definition~\ref{def:compat} requires: if $\forall
S.D \in \tau(e_i)$, then $D \in \tau'$.  So $D$ is already present
in~$\tau'$ at the moment of creation.

For part~(b): trivial, as $\forall S.D$ constrains only $S$-neighbors.

For the symmetric case: $R$-compatibility condition~(P2) requires that
if $\forall \inv(S_i).D \in \tau'$ and $(\tau(e_i), S_i, \tau') \in
P$, then $D \in \tau(e_i)$.  By \ref{inv:static}, $\tau(e_i)$ was
fixed at creation and already contains~$D$.
\end{proof}

\begin{remark}[Composition cannot force unexpected propagation]
\label{rem:composition-force}
One might worry that composition constraints could force an edge value
that triggers an unaccounted-for universal restriction.  For instance:
$R(e_a, e_b) = U$ and $R(e_b, e_{m+1}) = S_b$ might force $R(e_a,
e_{m+1}) = S_a$ where $S_a$ is the unique element of $\comp(U, S_b)$.
If $\forall S_a.D \in \tau(e_a)$ and $D \notin \tau'$, we would have
a problem.

This cannot happen.  The role $S_a$ was chosen from the domain $\{S
\mid (\tau(e_a), S, \tau') \in P\}$ by Claim~\ref{claim:extension}.
The patchwork property guarantees that a global assignment
exists within these domains while satisfying all composition
constraints.  Every role $S_a$ in the solution satisfies $(\tau(e_a),
S_a, \tau') \in P$, which by $R$-compatibility implies $D \in \tau'$
whenever $\forall S_a.D \in \tau(e_a)$.  The composition constraints
and the pair-compatibility constraints are solved \emph{jointly}; the
patchwork property ensures a simultaneous solution exists.
\end{remark}

\subsubsection*{Concept truth in the limit model}

We now prove that types faithfully determine concept truth.

\begin{lemma}[Concept truth]
\label{lem:truth}
In the limit model $\interp = \bigcup_{n} (\Delta_n,
\cdot^{\interp}_n)$, for every $e \in \DeltaI$ and every $D \in
\cl(C_0)$:
\[
  D \in \tau(e) \;\implies\; e \in D^{\interp}.
\]
\end{lemma}

\begin{proof}
By structural induction on~$D$.

\smallskip\noindent\textbf{Case $D = A$ (concept name):}\;
By construction, $A^{\interp} = \{e \in \DeltaI \mid A \in \tau(e)\}$.
If $A \in \tau(e)$, then $e \in A^{\interp}$.

\smallskip\noindent\textbf{Case $D = \neg A$:}\;
If $\neg A \in \tau(e)$, then $A \notin \tau(e)$ (by type condition
(T3): a type cannot contain both $A$ and $\neg A$).  So $e \notin
A^{\interp}$, hence $e \in (\neg A)^{\interp}$.

\smallskip\noindent\textbf{Case $D = D_1 \sqcap D_2$:}\;
If $D_1 \sqcap D_2 \in \tau(e)$, then $D_1 \in \tau(e)$ and $D_2 \in
\tau(e)$ by type condition~(T1).  By induction, $e \in
D_1^{\interp}$ and $e \in D_2^{\interp}$, so $e \in (D_1 \sqcap
D_2)^{\interp}$.

\smallskip\noindent\textbf{Case $D = D_1 \sqcup D_2$:}\;
If $D_1 \sqcup D_2 \in \tau(e)$, then $D_1 \in \tau(e)$ or $D_2 \in
\tau(e)$ by type condition~(T2).  By induction, $e \in
D_1^{\interp}$ or $e \in D_2^{\interp}$, so $e \in (D_1 \sqcup
D_2)^{\interp}$.

\smallskip\noindent\textbf{Case $D = \forall R.E$:}\;
Suppose $\forall R.E \in \tau(e)$ and $(e, e') \in R^{\interp}$
(i.e., $R(e, e') = R$ for some $e' \in \DeltaI$).  By
invariant~\ref{inv:pair}, $(\tau(e), R, \tau(e')) \in P$.  By
$R$-compatibility~(P1), $E \in \tau(e')$.  By induction hypothesis,
$e' \in E^{\interp}$.

Since this holds for \emph{every} $R$-neighbor $e'$, we have $e \in
(\forall R.E)^{\interp}$.

\emph{This is the step where dormant restrictions are handled.}  At
the stage when $e'$ is created, the edge $R(e, e')$ triggers the
universal restriction $\forall R.E \in \tau(e)$.  But
Lemma~\ref{lem:dormant} guarantees $E \in \tau(e')$ --- the concept
was placed in $\tau(e')$ at creation, not propagated after the fact.
No concept is added to $\tau(e)$ or $\tau(e')$ at any later stage.

\smallskip\noindent\textbf{Case $D = \exists R.E$:}\;
Suppose $\exists R.E \in \tau(e)$.  By the Henkin construction with
fair enumeration, there eventually exists a stage~$n$ at which the
demand $(e, \exists R.E)$ is processed: a new element $e'$ is created
with $\tau(e') = \tau'$ where $E \in \tau'$ and $R(e, e') = R$.  By
induction hypothesis, $e' \in E^{\interp}$.  Hence $e \in (\exists
R.E)^{\interp}$.
\end{proof}

\subsubsection*{The role of monotonicity along PP-chains}

The model constructed above may contain infinite $\PP$-chains (e.g.,
when $(\exists \PP.\top) \sqcap (\forall \PP.\exists \PP.\top) \in
\tau_0$).  The monotonicity lemmas (Section~\ref{sec:monotonicity})
explain why the finite quasimodel faithfully represents these infinite
structures.

Consider an infinite $\PP$-chain $e_0, e_1, e_2, \ldots$ in the
limit model (with $R(e_i, e_{i+1}) = \PP$ for all~$i$) and an
external element $x$.  By Lemma~\ref{lem:upward}, the relation $R(x,
e_i)$ follows a monotone progression:
\[
  \DR \;\to\; \DR, \PO, \PP; \quad
  \PO \;\to\; \PO, \PP; \quad
  \PP \;\to\; \PP; \quad
  \PPI \;\to\; \PO, \EQ, \PP, \PPI.
\]
After at most 3 transitions (e.g., $\DR \to \PO \to \PP$), the
relation \emph{stabilizes}: once $R(x, e_k) = \PP$ for some~$k$,
$R(x, e_i) = \PP$ for all $i \geq k$.  By Lemma~\ref{lem:downward},
$\DR$ and $\PPI$ are downward persistent: once $R(x, e_k) = \DR$,
$R(x, e_i) = \DR$ for all $i \leq k$.

This has three consequences for the Henkin construction:

\begin{enumerate}[nosep]
  \item \textbf{Bounded type variation.}\; Along the chain, only
    finitely many distinct pair-types $(\tau(x), S, \tau(e_i))$ arise
    (since the relation $S = R(x, e_i)$ takes at most $|\NR
    \setminus \{\EQ\}| = 4$ values and stabilizes).  All such
    pair-types are in~$P$.

  \item \textbf{Uniform universal propagation.}\; Once $R(x, e_i)$
    stabilizes to some role~$S$, the same set of concepts $\{E \mid
    \forall S.E \in \tau(x)\}$ applies to all subsequent chain
    elements.  No \emph{new} concept enters any $\tau(e_i)$ from $x$'s
    perspective beyond the stabilization point.

  \item \textbf{Finite representability.}\; The types $\tau(e_i)$
    along the chain are drawn from~$T$ (finitely many).  The
    monotonicity ensures that the ``relational profile'' of the chain
    (the mapping from external elements to sequences of pair-types)
    eventually becomes periodic.  The quasimodel captures all relevant
    type/pair-type/triple-type information, so the infinite chain is
    fully determined by its finite abstraction.
\end{enumerate}

Without monotonicity, the relations along a $\PP$-chain could
potentially oscillate (e.g., $\DR \to \PO \to \DR \to \PO \to
\cdots$), generating an unbounded variety of pair-type configurations.
The monotonicity lemmas rule this out: the progression is
one-directional, ensuring that the finite quasimodel is a faithful
representation of the infinite model.

\subsubsection*{Verification of model properties}

The limit model $\interp = \bigcup_{n \geq 0} (\Delta_n,
\cdot^{\interp}_n)$ is a valid $\ALCIRCC{5}$ interpretation:
\begin{itemize}[nosep]
  \item \emph{One cluster:} every pair of elements has a role
    assigned (by construction, each new element receives edges to all
    existing elements).
  \item \emph{JEPD:} each pair has exactly one base relation (roles
    are assigned uniquely and permanently).
  \item \emph{Composition:} every triple satisfies the composition
    table (invariant~\ref{inv:comp}, maintained at each stage).
  \item \emph{Converse:} maintained by setting $R(e_{m+1}, e_i) =
    \inv(S_i)$.
  \item \emph{Strong EQ:} $R(e, e) = \EQ$ for all $e$, and no
    distinct pair has $\EQ$ (since $S_i \in \NR \setminus \{\EQ\}$).
  \item \emph{Concept truth:} by Lemma~\ref{lem:truth}, $C_0 \in
    \tau_0 = \tau(e_1)$ implies $e_1 \in C_0^{\interp}$.
\end{itemize}

\noindent
Therefore $\interp \models C_0$, and $C_0$ is satisfiable.
\end{proof}

%% ============================================================
\section{Decidability and Complexity via Type Elimination}
\label{sec:complexity}

\begin{theorem}[\textbf{RETRACTED}]
\label{thm:complexity}
\sout{Concept satisfiability in $\ALCIRCC{5}$ is decidable in
\EXPTIME.}
\end{theorem}

\noindent
\textbf{This theorem is wrong.}  The type-elimination algorithm
described below rejects satisfiable concepts.  See
Remark~\ref{rem:type-elim-bug} for the concrete counterexample and
root cause analysis.

\medskip
\noindent
\emph{The original proof is retained (struck through) for reference:}

\begin{proof}[\textbf{Original proof (incorrect)}]
By the soundness direction of Theorem~\ref{thm:main}, if $C_0$ is
satisfiable then a quasimodel for $C_0$ exists.  Equivalently, if no
quasimodel exists then $C_0$ is unsatisfiable.  We describe a
type-elimination algorithm that decides quasimodel existence.

\begin{enumerate}[nosep]
  \item \textbf{Initialization.}\; Compute $\Tp(C_0)$, $\Pair(C_0)$,
    and $\Trip(C_0)$.  These sets have sizes bounded by
    $2^{O(|C_0|)}$, $2^{O(|C_0|)}$, and $2^{O(|C_0|)}$ respectively.
    Set $T := \Tp(C_0)$ and $P := \{(\tau_1, R, \tau_2) \in \Pair(C_0)
    \mid \tau_1, \tau_2 \in T\}$.

  \item \textbf{Algebraic closure enforcement.}\; Enforce condition
    (Q3) on the type-level disjunctive network
    $\mathsf{DN}$: for each triple $(\tau_1, \tau_2, \tau_3)$ and each
    $R_{12} \in \mathsf{DN}(\tau_1, \tau_2)$, check whether there exist
    $R_{13} \in \mathsf{DN}(\tau_1, \tau_3)$ and $R_{23} \in
    \mathsf{DN}(\tau_2, \tau_3)$ with $R_{13} \in \comp(R_{12}, R_{23})$.
    If not, remove $(\tau_1, R_{12}, \tau_2)$ from~$P$ and repeat
    until stable.

  \item \textbf{Type elimination.}\; Iteratively remove from $T$ any
    type $\tau$ that fails one of:
    \begin{itemize}[nosep]
      \item \emph{Existential demand (Q1):} there exists $\exists R.D
        \in \tau$ for which no witness type $\tau' \in T$ with $D \in
        \tau'$ and $R \in \mathsf{DN}(\tau, \tau')$ can be found.
      \item \emph{Completeness (Q2):} there exists a distinct type
        $\tau' \in T$ with $\mathsf{DN}(\tau, \tau') = \emptyset$.
    \end{itemize}
    When a type is removed, also remove all pair-types involving it,
    re-run algebraic closure enforcement (Step~2), then repeat.

  \item \textbf{Termination.}\; Each iteration removes at least one
    type or pair-type.  Since there are at most $2^{O(|C_0|)}$ types
    and $2^{O(|C_0|)}$ pair-types, termination occurs after at most
    $2^{O(|C_0|)}$ iterations.

  \item \textbf{Acceptance.}\; Accept if and only if there remains a
    type $\tau_0$ with $C_0 \in \tau_0$ such that $\tau_0$ and the
    surviving types and pair-types form a quasimodel for~$C_0$
    (i.e., conditions (Q1)--(Q3) are satisfied).
\end{enumerate}

Each iteration inspects at most $2^{O(|C_0|)}$ types, pairs, and
triples.  Total time: $2^{O(|C_0|)} \cdot 2^{O(|C_0|)} =
2^{O(|C_0|)}$, which is in \EXPTIME.

\medskip\noindent\textbf{Correctness (INCORRECT).}\;
\sout{If $C_0$ is satisfiable, the soundness proof
(Section~\ref{sec:soundness}) produces a quasimodel $(T^*, P^*,
\tau_0)$ whose types, pair-types, and algebraic closure property are
all preserved by the elimination algorithm.  Hence $\tau_0$ survives,
and the algorithm accepts.}

\medskip\noindent
\textbf{The flaw:} The quasimodel $(T^*, P^*)$ extracted from a
model is \emph{not} preserved by the elimination algorithm.  The
algorithm starts with $T = \Tp(C_0) \supsetneq T^*$, and (Q3)
enforcement with the full type set~$T$ removes pair-types that are
valid within~$T^*$.  This is because (Q3) universally quantifies
over \emph{all} $\tau_3 \in T$: types in $T \setminus T^*$ may have
$\mathsf{DN}(\tau_2, \tau_3) = \emptyset$, causing (Q3) to remove
$(\tau_1, R_{12}, \tau_2)$ even though the triple $(\tau_1, \tau_2,
\tau_3)$ never arises in the model.  The resulting (Q1) and (Q2)
cascade eliminates all types, including $\tau_0$.  See
Remark~\ref{rem:type-elim-bug}.
\end{proof}

\begin{remark}[The type elimination bug]
\label{rem:type-elim-bug}
The correctness argument for Theorem~\ref{thm:complexity} fails because
the greatest-fixpoint elimination does not preserve model-extracted
quasimodels.  The flaw is subtle: condition~(Q3) is
\emph{anti-monotone} in the type set~$T$ (larger $T$ makes (Q3)
harder to satisfy), while~(Q1) is anti-monotone in the opposite
direction (smaller $T$ makes (Q1) harder).  The greatest-fixpoint
approach assumes these conditions are compatible, but they are not.

\textbf{Concrete counterexample.}\;
$C_1 = \exists \PO.D \sqcap \exists \DR.(B \sqcap \forall \PO.\neg D)
\sqcap \forall \DR.\neg D \sqcap \forall \PP.\neg D \sqcap \forall
\PPI.\neg D$
is satisfiable: the model $\{e_0, e_1, e_2\}$ with $e_0$~(root,
$\neg B, \neg D$), $e_1$~($B$-filler, $B, \neg D$), $e_2$~($D$-element,
$D, \neg B$), and edges $R(e_0, e_1) = \DR$, $R(e_0, e_2) = \PO$,
$R(e_1, e_2) = \DR$ satisfies $C_1$ at $e_0$ (composition-consistent,
verified computationally).

The three model-extracted types $\{\tau_0, \tau_1, \tau_2\}$ with
compatible DN $=$ $\{(\tau_0, \tau_1): \{\DR, \PO, \PPI\},\;
(\tau_0, \tau_2): \{\PO\},\; (\tau_1, \tau_2): \{\DR\}\}$ form a
valid quasimodel: (Q1)~$\checkmark$, (Q2)~$\checkmark$,
(Q3)~$\checkmark$ (pairwise-distinct triples).

However, $|\Tp(C_1)| = 128$.  Among these, types with~$D$ and
$\forall R.\neg D$ for all~$R$ have empty DN with all other $D$-types.
The elimination cascade:
\begin{enumerate}[nosep]
  \item (Q3) with all 128 types prunes DN entries: $(\tau_0, \PP,
    \tau_1)$ is removed because there exists $\tau_3$ (a $D$-type
    in $\Tp(C_1)$ with empty DN to $\tau_1$).
  \item (Q2) kills types whose DN to some surviving type becomes
    empty.  This removes 78 types in the first iteration.
  \item (Q1) kills types that lost their witnesses: 50~more types
    in the second iteration.
  \item After 2 iterations: $T = \emptyset$, 0~types survive.
\end{enumerate}

\textbf{The fundamental issue.}\;
For standard $\ALCI$ (no spatial constraints), type elimination uses
only~(Q1).  (Q1) is anti-monotone: the greatest fixpoint correctly
finds the largest self-sustaining type set.  For $\ALCIRCC{5}$,
(Q3) adds a condition that is anti-monotone in the \emph{opposite}
direction.  These opposing forces make the greatest-fixpoint approach
incorrect.  A correct algorithm would need \emph{subset search}
(find $S \subseteq \Tp(C_0)$ where (Q1)--(Q3) hold internally),
which is computationally harder.
\end{remark}

\begin{theorem}[Lower bound]
\label{thm:lower}
Concept satisfiability in $\ALCIRCC{5}$ is \PSPACE-hard.
\end{theorem}

\begin{proof}
Wessel~\cite[Section 6.3]{Wessel2003} establishes this by a reduction
from QBF validity, using the $\PO$-free fragment to enforce
tree-shaped models.
\end{proof}

\begin{corollary}
Concept satisfiability in $\ALCIRCC{5}$ is between \PSPACE{} and
\EXPTIME.
\end{corollary}

\begin{remark}[Closing the gap]
We conjecture that $\ALCIRCC{5}$ is \EXPTIME-complete.  The
interaction between $\PP$-transitivity and the universal role (which
enables TBox internalization) plausibly supports a reduction from the
\EXPTIME-hard \emph{two-player corridor tiling problem}, but we have
not established this.
\end{remark}


%% ============================================================
%% ============================================================
\section{A Tableau Decision Procedure}
\label{sec:tableau}

We now present an alternative decision procedure in the form of a
\emph{tableau calculus with blocking}.  Unlike the type-elimination
algorithm, which works globally on the space of all types and
pair-types, the tableau constructs a model-like structure
incrementally, expanding individual nodes.  The key challenge is that
$\ALCIRCC{5}$ models are complete graphs, not trees: creating a new
node requires assigning it a base relation to \emph{every} existing
node.  The patchwork property (Theorem~\ref{thm:patchwork}) is what
makes this tractable.

\subsection{Completion Graphs}

\begin{definition}[Completion graph]
\label{def:completion-graph}
A \emph{completion graph} for $C_0$ is a tuple $\mathcal{G} = (V,
\mathcal{L}, \mathcal{E}, \prec)$ where:
\begin{itemize}[nosep]
  \item $V$ is a finite, non-empty set of \emph{nodes}.
  \item $\mathcal{L} \colon V \to 2^{\cl(C_0)}$ is a \emph{labelling
    function} assigning a set of concepts to each node.
  \item $\mathcal{E} \colon \{(x,y) \in V \times V \mid x \neq y\}
    \to \NR \setminus \{\EQ\}$ is an \emph{edge labelling function}
    assigning a base relation to every ordered pair of distinct nodes,
    respecting converse: $\mathcal{E}(y,x) = \inv(\mathcal{E}(x,y))$.
  \item $\prec$ is a strict total order on $V$ (the \emph{creation
    order}).
\end{itemize}
The \emph{initial completion graph} for $C_0$ consists of a single
node $x_0$ with $\mathcal{L}(x_0) = \{C_0\}$, $V = \{x_0\}$, and
$\mathcal{E}$ vacuously defined (no pairs of distinct nodes yet).
\end{definition}

\begin{remark}
The edge labelling $\mathcal{E}$ is a \emph{total} function on
distinct pairs, reflecting the one-cluster and JEPD requirements.
Whenever a new node $y$ is added to $V$, the function $\mathcal{E}$
must be extended to assign a base relation $\mathcal{E}(z,y)$ for
every existing node $z$.  This is the principal difference from
tree-based tableau calculi, where a new node has a single edge to its
parent.
\end{remark}

\begin{definition}[Blocking]
\label{def:blocking}
Let $\mathcal{G} = (V, \mathcal{L}, \mathcal{E}, \prec)$ be a
completion graph.  A node $x \in V$ is \emph{blocked} if there exists
$y \prec x$ such that $y$ is not blocked and $\mathcal{L}(x) =
\mathcal{L}(y)$.  We call $y$ a \emph{blocker} for $x$.  A node is
\emph{active} if it is not blocked.
\end{definition}

\begin{remark}
This is \emph{equality anywhere-blocking}: the blocker $y$ need not be
a predecessor or neighbor in any structural sense---it need only
precede $x$ in the creation order and carry the same label.  The
simplicity of this blocking condition (compared to, e.g., the pairwise
blocking of $\ALCI$ tree tableaux) is made possible by the patchwork
property, which allows us to decouple the type-level structure from
the specific relational neighborhood (see the soundness proof,
Section~\ref{sec:tab-soundness}).
\end{remark}

\subsection{Expansion Rules}
\label{sec:tab-rules}

The tableau calculus consists of four deterministic rules and two
nondeterministic rules, applied to a completion graph $\mathcal{G}$.
We write $\mathcal{L}(x)$ for the label and $\mathcal{E}(x,y)$ for
the edge label.

\medskip
\noindent\textbf{Propositional rules} (deterministic):

\smallskip
\noindent$(\sqcap)$\quad If $D_1 \sqcap D_2 \in \mathcal{L}(x)$ and
$\{D_1, D_2\} \not\subseteq \mathcal{L}(x)$: \\
\hspace*{2em} add $D_1$ and $D_2$ to $\mathcal{L}(x)$.

\smallskip
\noindent$(\sqcup)$\quad If $D_1 \sqcup D_2 \in \mathcal{L}(x)$ and
$D_1 \notin \mathcal{L}(x)$ and $D_2 \notin \mathcal{L}(x)$: \\
\hspace*{2em} \emph{nondeterministically} add $D_1$ or $D_2$ to
$\mathcal{L}(x)$.

\medskip
\noindent\textbf{Propagation rule} (deterministic):

\smallskip
\noindent$(\forall)$\quad If $\forall R.D \in \mathcal{L}(x)$ and
$\mathcal{E}(x,y) = R$ and $D \notin \mathcal{L}(y)$: \\
\hspace*{2em} add $D$ to $\mathcal{L}(y)$.

\smallskip
\noindent This rule also handles inverses: since $\forall R^-.D$ is
equivalent to $\forall \inv(R).D$ in NNF, and $\mathcal{E}$ stores
base relations, the rule fires whenever $\forall S.D \in
\mathcal{L}(x)$ and $\mathcal{E}(x,y) = S$.

\medskip
\noindent\textbf{Generating rule} (nondeterministic):

\smallskip
\noindent$(\exists)$\quad If $\exists R.D \in \mathcal{L}(x)$, $x$ is
active, and there is no $y \in V$ with $\mathcal{E}(x,y) = R$ and
$D \in \mathcal{L}(y)$:
\begin{enumerate}[nosep,label=\arabic*.]
  \item Create a fresh node $y$ and add it to $V$ as the
    $\prec$-greatest element.
  \item Set $\mathcal{L}(y) := \{D\}$.
  \item Set $\mathcal{E}(x,y) := R$ and $\mathcal{E}(y,x) := \inv(R)$.
  \item For each $z \in V \setminus \{x, y\}$,
    \emph{nondeterministically} choose $S_z \in \NR \setminus \{\EQ\}$
    and set $\mathcal{E}(z,y) := S_z$, $\mathcal{E}(y,z) :=
    \inv(S_z)$.
\end{enumerate}

\begin{remark}
Step~4 is the hallmark of the complete-graph tableau: the new node
$y$ receives an edge to \emph{every} existing node.  The choice of
$S_z$ is nondeterministic, subject to the constraint-filtering
condition below.  The patchwork property guarantees that a valid
choice exists whenever the concept is satisfiable (see
Section~\ref{sec:tab-completeness}).
\end{remark}

\subsection{Constraint Filtering}
\label{sec:constraint-filtering}

After the $(\exists)$-rule creates a new node~$y$ with edges to all
existing nodes, we require:

\medskip\noindent\textbf{(CF) Composition consistency.}\; For every
triple of distinct nodes $(u, v, w)$ that includes $y$:
\[
  \mathcal{E}(u,w) \in \comp(\mathcal{E}(u,v),\, \mathcal{E}(v,w)).
\]

If no assignment of the $S_z$ values in step~4 of $(\exists)$
satisfies (CF) for all such triples simultaneously, the current branch
\emph{closes} (clash by composition failure).

\begin{remark}
Checking~(CF) amounts to solving an atomic RCC5 constraint network:
the existing nodes have fixed pairwise relations, and the new node~$y$
has the relation to~$x$ fixed at~$R$, with the remaining relations
$S_z$ as variables.  By Theorem~\ref{thm:patchwork}, a simultaneous
satisfying assignment exists if and only if every triple involving~$y$
is individually satisfiable, i.e., the network is path-consistent.
Path consistency can be checked (and enforced) in time $O(|V|^3)$.
\end{remark}

\subsection{Rule Application Strategy}
\label{sec:strategy}

To ensure that labels stabilize before blocking is evaluated:

\begin{enumerate}[nosep]
  \item Apply $(\sqcap)$, $(\sqcup)$, and $(\forall)$ exhaustively
    to all nodes (in any order, resolving $(\sqcup)$
    nondeterministically).
  \item Re-evaluate blocking: a node $x$ is blocked if
    $\mathcal{L}(x) = \mathcal{L}(y)$ for some non-blocked $y \prec x$.
  \item Apply one instance of $(\exists)$ for an active node with an
    unsatisfied demand.
  \item After adding the new node and its edges, apply $(\forall)$
    exhaustively (the new edges may trigger propagations both to and
    from the new node; these may cascade to other nodes).
  \item Re-evaluate blocking (label changes from $(\forall)$ may
    affect blocking status).
  \item Repeat from step~1 until no rule applies or a clash is
    detected.
\end{enumerate}

\begin{definition}[Saturated completion graph]
\label{def:saturated}
A completion graph $\mathcal{G}$ is \emph{saturated} if no expansion
rule is applicable (i.e., all propositional rules have been applied,
all $(\forall)$-propagations have been performed, and every active
node's existential demands are witnessed by an appropriate neighbor).
\end{definition}

\begin{definition}[Clash-free completion graph]
\label{def:clash-free}
A completion graph is \emph{clash-free} if:
\begin{enumerate}[label=(C\arabic*),nosep]
  \item No node has $\{A, \neg A\} \subseteq \mathcal{L}(x)$ for any
    concept name $A$.
  \item Composition consistency (CF) holds for all triples.
\end{enumerate}
A completion graph that is both saturated and clash-free is called
\emph{open}.
\end{definition}

\noindent
The tableau \emph{decides} satisfiability: $C_0$ is satisfiable iff
there exists a sequence of nondeterministic choices producing an open
completion graph.

\subsection{Termination}
\label{sec:tab-termination}

\begin{theorem}[Termination]
\label{thm:termination}
Every sequence of rule applications terminates.  The final completion
graph has at most $2^{O(|C_0|)}$ nodes.
\end{theorem}

\begin{proof}
Let $n = |\cl(C_0)| \leq 2|C_0|$ and $N = 2^n = |\Tp(C_0)|$.

\smallskip\noindent\textbf{Bound on active nodes.}\; At any point,
the active (non-blocked) nodes have pairwise distinct labels (by the
definition of blocking).  Each label is a subset of $\cl(C_0)$, so
there are at most $N$ active nodes.

\smallskip\noindent\textbf{Labels grow monotonically.}\; The rules
$(\sqcap)$, $(\sqcup)$, and $(\forall)$ only \emph{add} concepts to
labels, never remove them.  Therefore, each node's label changes at
most $n$ times (once per concept in $\cl(C_0)$).

\smallskip\noindent\textbf{Bound on node creation.}\; The $(\exists)$
rule fires only for active nodes with unsatisfied demands.  Each
active node~$x$ has at most $n$ existential demands (concepts of the
form $\exists R.D$ in $\cl(C_0)$).  Each demand generates at most one
new node (once a witness exists, the demand is satisfied and the rule
no longer fires for that demand).  With at most $N$ active nodes, at
most $N \cdot n$ nodes are created by the $(\exists)$ rule.  Each
newly created node is either active (contributing to the at-most-$N$
budget) or immediately blocked.

\smallskip\noindent\textbf{Bound on re-evaluations.}\; When a label
changes (via $(\forall)$ propagation), blocking status may change, and
a previously blocked node may become active, requiring expansion.
However, each label change is permanent (labels only grow), and the
total number of label changes across all nodes is at most $|V| \cdot
n$.  Each unblocking event is caused by a label change, so the total
number of unblocking events is bounded.  Each unblocking event may
generate at most $n$ new witnesses, but the total number of witnesses
is bounded by the total number of active-node $\times$ demand pairs
encountered during the process, which is at most $N \cdot n$ per
``generation'' of the label.

\smallskip\noindent\textbf{Overall bound.}\; The total number of
nodes is bounded by $N \cdot n \cdot n = 2^n \cdot n^2 =
2^{O(|C_0|)}$, and the total number of rule applications (label
additions and node creations) is bounded by $|V| \cdot n =
2^{O(|C_0|)}$.  Hence the procedure terminates in at most
$2^{O(|C_0|)}$ steps.
\end{proof}

\subsection{Soundness: Open Branch $\Rightarrow$ Satisfiable}
\label{sec:tab-soundness}

\begin{theorem}[Soundness of the tableau]
\label{thm:tab-soundness}
If the tableau produces an open \textup{(}saturated, clash-free\textup{)}
completion graph $\mathcal{G}$ for $C_0$, then $C_0$ is satisfiable
in $\ALCIRCC{5}$.
\end{theorem}

\begin{proof}
We show that an open completion graph $\mathcal{G} = (V, \mathcal{L},
\mathcal{E}, \prec)$ encodes a valid quasimodel.  The result then
follows from Theorem~\ref{thm:main}.

\medskip\noindent\textbf{Step 1: Quasimodel extraction.}\; Define:
\begin{align*}
  T &= \{ \mathcal{L}(x) \mid x \in V,\; x \text{ is active} \}, \\
  P &= \{ (\mathcal{L}(x), \mathcal{E}(x,y), \mathcal{L}(y))
    \mid x,y \in V,\; x \neq y \}, \\
  \tau_0 &= \mathcal{L}(x_0).
\end{align*}

\medskip\noindent\textbf{Step 2: Verification of (Q1).}\; Let $\tau
\in T$ and $\exists R.D \in \tau$.  Then $\tau = \mathcal{L}(x)$ for
some active node $x$.  Since $\mathcal{G}$ is saturated, there exists
$y \in V$ with $\mathcal{E}(x,y) = R$ and $D \in \mathcal{L}(y)$.
Set $\tau' = \mathcal{L}(y)$.  Then $D \in \tau'$ and
$R = \mathcal{E}(x,y) \in \mathsf{DN}(\tau, \tau')$, verifying~(Q1).

\medskip\noindent\textbf{Step 3: Verification of (Q2).}\; Since
$\mathcal{E}$ is a total function on distinct pairs, for any distinct
$\tau_1, \tau_2 \in T$ with active witnesses $x_1, x_2$, the edge
$\mathcal{E}(x_1, x_2)$ provides a role in
$\mathsf{DN}(\tau_1, \tau_2)$, so
$\mathsf{DN}(\tau_1, \tau_2) \neq \emptyset$.

\medskip\noindent\textbf{Step 4: Verification of (Q3).}\;
Let $\tau_1, \tau_2, \tau_3 \in T$ (not necessarily distinct) and
$R_{12} \in \mathsf{DN}(\tau_1, \tau_2)$.  Then there exist nodes
$a, b \in V$ with $\mathcal{L}(a) = \tau_1$, $\mathcal{L}(b) =
\tau_2$, and $\mathcal{E}(a,b) = R_{12}$.  Since $\tau_3 \in T$,
there exists a node $c \in V$ with $\mathcal{L}(c) = \tau_3$
(choosing any witness of $\tau_3$; if $\tau_3 = \tau_1$ we may take
$c = a$, and similarly for $\tau_3 = \tau_2$; for the self-pair case
$\tau_1 = \tau_3$ with $c = a$, we have $R_{13} = \EQ$ and
$\comp(R_{12}, R_{23}) \ni R_{12}$ trivially).

For the case of three distinct nodes $a, b, c$: the completion graph
assigns $\mathcal{E}(a,c)$ and $\mathcal{E}(b,c)$.  The
composition-consistency condition~(CF) on the triple $(a, b, c)$ gives
\[
  \mathcal{E}(a,c) \in \comp(\mathcal{E}(a,b), \mathcal{E}(b,c))
  = \comp(R_{12}, \mathcal{E}(b,c)).
\]
Set $R_{13} = \mathcal{E}(a,c) \in \mathsf{DN}(\tau_1, \tau_3)$
and $R_{23} = \mathcal{E}(b,c) \in \mathsf{DN}(\tau_2, \tau_3)$.
Then $R_{13} \in \comp(R_{12}, R_{23})$, which is exactly the
algebraic closure condition~(Q3).

\medskip\noindent\textbf{Step 5: Conclusion.}\; The extracted triple
$(T, P, \tau_0)$ is a quasimodel for $C_0$.  Moreover, this
quasimodel arises from a model-like structure (the completion graph)
in which all pair-types are realized by concrete node pairs.  By the
completeness argument of Theorem~\ref{thm:main}, and the
non-emptiness guarantee for model-derived quasimodels
(Remark~\ref{rem:extension-gap}), $C_0$ is satisfiable in
$\ALCIRCC{5}$.
\end{proof}

\subsection{Completeness: Satisfiable $\Rightarrow$ Open Branch}
\label{sec:tab-completeness}

\begin{theorem}[Completeness of the tableau]
\label{thm:tab-completeness}
If $C_0$ is satisfiable in $\ALCIRCC{5}$, then there exists a
sequence of nondeterministic choices yielding an open completion graph
for $C_0$.
\end{theorem}

\begin{proof}
Let $\interp = (\DeltaI, \cdot^{\interp})$ be a model with $d_0 \in
C_0^{\interp}$.  We show that nondeterministic choices can be guided
by~$\interp$ so that no clash arises and every demand is satisfied.

\medskip\noindent\textbf{Guide function.}\;  We maintain a partial
function $\pi \colon V \to \DeltaI$ that maps each tableau node to a
``corresponding'' element in the model, with the invariant:
\begin{enumerate}[label=(\Roman*),nosep]
  \item $\mathcal{L}(x) \subseteq \tau^{\interp}(\pi(x))$ for all
    $x \in V$.
  \item $\mathcal{E}(x,y) = R^{\interp}(\pi(x), \pi(y))$ for all
    distinct $x, y \in V$ (where $R^{\interp}(a,b)$ denotes the
    unique base relation holding between $a$ and $b$ in $\interp$).
\end{enumerate}

\medskip\noindent\textbf{Initialization.}\;
Set $\pi(x_0) = d_0$.  Then $\mathcal{L}(x_0) = \{C_0\}$ and $d_0
\in C_0^{\interp}$, so invariant~(I) holds.

\medskip\noindent\textbf{Guiding the $(\sqcup)$ rule.}\;
When $D_1 \sqcup D_2 \in \mathcal{L}(x)$, invariant~(I) gives $D_1
\sqcup D_2 \in \tau^{\interp}(\pi(x))$, so $\pi(x) \in (D_1 \sqcup
D_2)^{\interp}$, meaning $\pi(x) \in D_1^{\interp}$ or $\pi(x) \in
D_2^{\interp}$.  Choose the true disjunct.  Invariant~(I) is
maintained.

\medskip\noindent\textbf{Guiding the $(\exists)$ rule.}\;
When $\exists R.D \in \mathcal{L}(x)$ and $x$ is active, invariant~(I)
gives $\pi(x) \in (\exists R.D)^{\interp}$, so there exists $d' \in
\DeltaI$ with $(\pi(x), d') \in R^{\interp}$ and $d' \in
D^{\interp}$.

Create node $y$ and set $\pi(y) = d'$.  Then:
\begin{itemize}[nosep]
  \item $\mathcal{E}(x,y) = R = R^{\interp}(\pi(x), \pi(y))$.
    Invariant~(II) holds for the pair $(x,y)$.
  \item For each existing node $z$, choose $S_z =
    R^{\interp}(\pi(z), \pi(y))$ (the actual relation in the model).
    Set $\mathcal{E}(z,y) = S_z$.  Invariant~(II) holds for $(z,y)$.
  \item Composition consistency~(CF) holds because the model satisfies
    the composition table.
\end{itemize}

\noindent After $(\forall)$-propagation, every concept added to
$\mathcal{L}(y)$ is of the form $D'$ where $\forall S.D' \in
\mathcal{L}(z)$ and $\mathcal{E}(z,y) = S$.  By invariant~(I),
$\forall S.D' \in \tau^{\interp}(\pi(z))$, so $(\pi(z), \pi(y)) \in
S^{\interp}$ implies $\pi(y) \in (D')^{\interp}$, i.e., $D' \in
\tau^{\interp}(\pi(y))$.  Similarly for propagation from $y$ back to
existing nodes.  Hence invariant~(I) is maintained for all affected
labels.

\medskip\noindent\textbf{Clash-freeness.}\;
Invariant~(I) implies $\mathcal{L}(x) \subseteq
\tau^{\interp}(\pi(x))$, which is a consistent type, so $\{A, \neg A\}
\not\subseteq \mathcal{L}(x)$.  Invariant~(II) implies
$\mathcal{E}(x,y) = R^{\interp}(\pi(x), \pi(y))$, and the model
satisfies composition, so~(CF) holds.

\medskip\noindent\textbf{Demand satisfaction despite blocking.}\;
The guide function shows that all nondeterministic choices can be made
consistently.  The only remaining concern is that a blocked node $x$
might have an unsatisfied demand.

Let $x$ be blocked by $y$ (so $\mathcal{L}(x) = \mathcal{L}(y)$ and
$y$ is active).  Suppose $\exists R.D \in \mathcal{L}(x)$.  Then
$\exists R.D \in \mathcal{L}(y)$, and since $y$ is active and the
graph is saturated, there exists a node $w$ with $\mathcal{E}(y,w) =
R$ and $D \in \mathcal{L}(w)$.

The demand of $x$ is $\exists R.D$.  We need an $R$-neighbor of $x$
containing $D$.  Node $w$ is an $\mathcal{E}(y,w)$-neighbor of $y$
with role $R$, but $\mathcal{E}(x,w)$ might be a different role.

This does not cause a problem \emph{for the tableau}: the tableau does
not require blocked nodes' demands to be directly satisfied in the
completion graph.  The completion graph is not itself a model; it is a
finite representation from which a model is constructed (via the
quasimodel extraction and Henkin construction in
Theorem~\ref{thm:tab-soundness}).  The quasimodel extracted in the
soundness proof has the type $\mathcal{L}(x) = \mathcal{L}(y) = \tau
\in T$, and the demand $\exists R.D$ of $\tau$ is witnessed by the
type $\mathcal{L}(w) = \tau'$ via the pair-type $(\tau, R, \tau') \in
P$.  The Henkin construction produces a model satisfying all demands of
all types, including those of blocked nodes.

\medskip\noindent The existence of the guided sequence of choices
shows that the nondeterministic tableau has at least one execution
producing a saturated, clash-free completion graph.
\end{proof}

\subsection{Summary of the Tableau Procedure}

\begin{corollary}[\textbf{NOT ESTABLISHED}]
\label{cor:tableau-decides}
\sout{The tableau calculus of Section~\ref{sec:tab-rules}, with blocking
\textup{(}Definition~\ref{def:blocking}\textup{)} and constraint
filtering \textup{(}Section~\ref{sec:constraint-filtering}\textup{)},
is a decision procedure for $\ALCIRCC{5}$ concept satisfiability.}
\end{corollary}

\noindent
\textbf{This corollary is not established.}  The individual components:
\begin{itemize}[nosep]
  \item \textbf{Termination}: established
    (Theorem~\ref{thm:termination}).
  \item \textbf{Completeness}: established
    (Theorem~\ref{thm:tab-completeness}).
  \item \textbf{Soundness}: \textbf{not established}.
    Theorem~\ref{thm:tab-soundness} extracts a quasimodel from an
    open completion graph and invokes the Henkin construction
    (Theorem~\ref{thm:main}), which has the extension gap
    (Remark~\ref{rem:extension-gap}).  The extension step requires
    solving a disjunctive constraint network at each stage; 1,911
    counterexamples to pairwise-implies-global solvability were found
    computationally at $m = 3$.  Without closing the extension gap,
    the tableau's soundness is unproven.
\end{itemize}

\begin{remark}[Comparison with tree-based tableaux]
\label{rem:comparison}
In standard $\ALCI$ tableaux, the completion structure is a forest,
new nodes have a single edge to their parent, and blocking uses
\emph{pairwise} matching (comparing both a node and its parent with
the blocker and its parent) to handle inverse roles.  In the
$\ALCIRCC{5}$ tableau:
\begin{enumerate}[nosep]
  \item The completion structure is a \emph{complete graph}, not a
    tree.
  \item The $(\exists)$-rule creates edges to \emph{all} existing
    nodes, not just one.
  \item Composition consistency~(CF) replaces the tree-structure
    invariant.
  \item Blocking is \emph{simpler}---just type equality, with no
    need for pairwise matching---because the patchwork property handles
    the relational context: the quasimodel extraction
    (Theorem~\ref{thm:tab-soundness}) shows that type-level
    information suffices to guarantee model existence, regardless of
    the specific role assignments between nodes.
  \item Soundness uses an indirect argument (quasimodel extraction +
    Henkin construction) rather than directly reading off the model
    from the completion graph.
\end{enumerate}
\end{remark}

\begin{remark}[Practical optimizations]
An implementation of this tableau can benefit from:
\begin{itemize}[nosep]
  \item \emph{Lazy role assignment}: Instead of choosing all
    $\mathcal{E}(z,y)$ at node creation, maintain a constraint network
    and only instantiate roles when needed for $(\forall)$-propagation.
  \item \emph{Path consistency preprocessing}: Before node creation,
    run arc consistency on the constraint network to prune impossible
    role values.
  \item \emph{Dependency-directed backtracking}: When a clash is
    detected, track which $(\sqcup)$ and $(\exists)$-choices
    contributed, and backtrack to the most recent relevant choice.
  \item \emph{Caching}: Since blocking is type-based, previously
    explored subtrees for a given type can be reused.
\end{itemize}
\end{remark}


%% ============================================================
\section{Extension to $\ALCIRCC{8}$}
\label{sec:rcc8}

The RCC8 calculus uses eight base relations: $\DC$, $\EC$, $\PO$,
$\EQ$, $\TPP$, $\TPPI$, $\NTPP$, $\NTPPI$.  The logic $\ALCIRCC{8}$
is defined analogously to $\ALCIRCC{5}$, using the RCC8 composition
table.  Both decision procedures extend to $\ALCIRCC{8}$, but the
transfer requires care because RCC8 and RCC5 differ in an important
tractability property.

\paragraph{What transfers directly.}
The \emph{atomic} patchwork property holds for RCC8: path-consistent
atomic RCC8 networks are globally consistent
(Theorem~\ref{thm:patchwork}).  Since the quasimodel and tableau
constructions rely on this atomic property, the following components
carry over unchanged:
\begin{itemize}[nosep]
  \item Types, pair-types, and the quasimodel conditions
    (Q1)--(Q3) are defined over the seven non-EQ RCC8 base relations
    $\{\DC, \EC, \PO, \TPP, \TPPI, \NTPP, \NTPPI\}$, using the RCC8
    composition table.
  \item The soundness proof (satisfiable concept $\Rightarrow$
    quasimodel) extracts types and pair-types from a model exactly as
    in Section~\ref{sec:soundness}.  Condition~(Q3) is verified from
    the model's composition closure.
  \item The type-elimination algorithm
    (Section~\ref{sec:complexity}) enforces algebraic closure and
    iteratively eliminates types; it runs in $2^{O(|C_0|)}$ time (the
    constant increases from 4 to 7 non-EQ relations, but the
    exponential bound is unchanged).
  \item The tableau calculus uses the RCC8 composition table for
    constraint filtering~(CF).  Blocking, termination, and the
    extraction of a quasimodel from an open completion graph proceed
    as in Section~\ref{sec:tableau}.
\end{itemize}

\paragraph{The tractability difference.}
For RCC5, the \emph{full} algebra is tractable: disjunctive
(non-atomic) path-consistent RCC5 networks are consistent
(Theorem~\ref{thm:rcc5-tract}).  For RCC8, this fails:
consistency of disjunctive RCC8 constraint networks is
\NP-complete~\cite{RenzNebel1999}.  The maximal tractable
subalgebras of RCC8 are proper subsets of the full
algebra~\cite{Renz1999}.

In the RCC5 proof, full tractability is invoked in
Claim~\ref{claim:extension} (Step~3) to solve the disjunctive
extension network that arises during the Henkin construction.  For
RCC8, this step is not available.  However, the argument for
\emph{model-derived} quasimodels does not depend on full
tractability: the model itself provides a globally consistent role
assignment that survives the path-consistency enforcement (see the
non-emptiness paragraph in the proof of Claim~\ref{claim:extension}).
Consequently:
\begin{itemize}[nosep]
  \item The soundness direction of the characterization theorem
    (satisfiable $\Rightarrow$ quasimodel) holds for RCC8.
  \item The completeness direction (quasimodel $\Rightarrow$ model
    via Henkin) holds for quasimodels extracted from models, by the
    same argument as for RCC5.
  \item The extension gap (Remark~\ref{rem:extension-gap}) applies
    equally: for \emph{abstract} quasimodels not arising from models,
    the Henkin construction may fail.  Unlike RCC5, where full
    tractability provides an alternative route to solving the extension
    network, no such fallback is available for RCC8.
  \item \textbf{The type-elimination algorithm is unsound for RCC8
    as well}, since the same anti-monotonicity of~(Q3) applies
    (see Remark~\ref{rem:type-elim-bug}).  The RCC8 algorithm
    inherits the cascade-elimination bug from the RCC5 version.
  \item The tableau's soundness proof extracts a quasimodel from an
    open completion graph---a structure where all pair-types are
    realized by concrete node pairs and composition consistency~(CF)
    holds---and invokes the Henkin construction for model-derived
    quasimodels.
\end{itemize}

\begin{theorem}[\textbf{RETRACTED}]
\label{thm:rcc8}
\sout{Concept satisfiability in $\ALCIRCC{8}$ is decidable in \EXPTIME.}
\end{theorem}

\noindent
\textbf{This theorem is wrong.}  It depends on the type-elimination
algorithm (Theorem~\ref{thm:complexity}), which is unsound.  The same
anti-monotonicity of~(Q3) that invalidates the RCC5 algorithm also
invalidates the RCC8 version.

\begin{proof}[\textbf{Original proof (incorrect)}]
By the analysis above: the type-elimination algorithm is a correct
\EXPTIME{} decision procedure for $\ALCIRCC{8}$.  The soundness
direction uses the RCC8 composition table and is independent of
tractability.  The algorithm's acceptance criterion (survival of
$\tau_0$ through elimination) ensures no false negatives.  The
\EXPTIME{} time bound follows from the same analysis as for
$\ALCIRCC{5}$ (Section~\ref{sec:complexity}), with the constant
factor increasing from $4$ to $7$ non-EQ base relations.
\end{proof}

\begin{theorem}[Wessel~\cite{Wessel2003}]
\label{thm:rcc8lower}
Concept satisfiability in $\ALCIRCC{8}$ is \EXPTIME-hard.
\end{theorem}

\begin{proof}
By reduction from the two-player corridor tiling game~\cite[Section
6.4]{Wessel2003}.
\end{proof}

\begin{corollary}[\textbf{RETRACTED}]
\label{cor:rcc8complete}
\sout{Concept satisfiability in $\ALCIRCC{8}$ is \EXPTIME-complete.}
\end{corollary}

\noindent
\textbf{Retracted:} depends on Theorem~\ref{thm:rcc8}, which is wrong.
The \EXPTIME{} lower bound (Theorem~\ref{thm:rcc8lower}) remains valid.
Decidability of $\ALCIRCC{8}$ is \textbf{open}.

%% ============================================================
\section{Related Work}
\label{sec:related}

\paragraph{Wessel's $\ALCRASG$ (2000/2001).}
The $\ALCRASG$ framework~\cite{Wessel2000} handles role boxes that
are \emph{admissible}: deterministic (each composition yields a single
relation), functional, complete, and associative.  For such role
boxes, the logic reduces to $\ALC$ with general TBoxes, yielding
\EXPTIME-completeness.  The RCC5 and RCC8 role boxes are
non-deterministic (many composition entries yield sets of relations)
and hence fall outside $\ALCRASG$.  Our quasimodel and tableau methods
handle this non-determinism by working with all possible completions
simultaneously, guided by the patchwork property.

\paragraph{Lutz \& Wolter (2006): Modal logics of topological relations.}
Lutz and Wolter~\cite{LutzWolter2006} study modal logics with RCC8
modalities interpreted over \emph{actual topological spaces}.  In
their setting, models must correspond to genuine spatial configurations
(e.g., in $\mathbb{R}^n$), which imposes additional constraints beyond
the composition table.  This leads to undecidability in many cases.
By contrast, $\ALCIRCC{}$ uses \emph{abstract} semantics: any
interpretation satisfying the composition table constraints is
admissible.  The abstract semantics are more permissive, and this
permissiveness is exactly what makes our decidability proof work.

\paragraph{Lutz \& Mili\v{c}i\'{c} (2007): $\omega$-admissible concrete domains.}
Lutz and Mili\v{c}i\'{c}~\cite{LutzMilicic2007} prove that $\ALC$
extended with $\omega$-admissible concrete domains is decidable.  RCC5
and RCC8 are $\omega$-admissible, so $\ALC(\mathrm{RCC5})$ and
$\ALC(\mathrm{RCC8})$ are decidable.  However, the concrete-domain
approach is fundamentally different: spatial predicates constrain
\emph{concrete features} of domain elements via functional roles, and
one cannot quantify \emph{over} spatial relations.  In particular,
$\forall \PP.C$ and $\exists \DR.D$ are not expressible in
$\ALC(\mathrm{RCC5})$.  The $\ALCIRCC{}$ logics are strictly more
expressive for spatial reasoning because spatial relations serve
directly as roles.

\paragraph{Borgwardt, De~Bortoli \& Koopmann (2024).}
Borgwardt et al.~\cite{BorgwardtEtAl2024} sharpen the complexity of
reasoning in $\ALC$ with $\omega$-admissible concrete domains.  Their
results apply to the concrete-domain setting and do not directly bear
on $\ALCIRCC{}$, but confirm that the patchwork/amalgamation property
of RCC calculi is a powerful tool for decidability in description
logics with spatial features.

\paragraph{Baader \& Rydval (2020): Concrete domains revisited.}
Baader and Rydval~\cite{BaaderRydval2020} strengthened undecidability
results for description logics with concrete domains and general
concept inclusions, and generalized the $\omega$-admissibility
conditions.  Their results refine the boundary between decidable and
undecidable concrete-domain extensions, but remain within the
concrete-domain paradigm (functional roles to concrete values) and do
not address composition-based role boxes.

\paragraph{Demri \& Gu (2026): Constraint automata with inverse roles.}
Demri and Gu~\cite{DemriGu2026} extend the automata-based approach to
description logics with concrete domains to handle \emph{inverse
roles}, functional role names, and constraint assertions, establishing
\EXPTIME{} membership.  This is the closest published result to
$\ALCIRCC{}$, as it combines inverse roles with RCC-like spatial
constraints.  However, it still operates within the concrete-domain
formalism: spatial predicates constrain concrete features via
functional roles, and one cannot quantify over spatial relations as
roles.  In $\ALCIRCC{5}$, concepts like $\forall\PP.C$ and
$\exists\DR.D$---which use spatial relations directly as
roles---remain inexpressible in the concrete-domain setting.  The two
formalisms are thus complementary: the concrete-domain approach is
well-suited for reasoning about spatial \emph{attributes} of domain
elements, while the $\ALCIRCC{}$ approach captures direct spatial
\emph{relationships} between elements.

\paragraph{Constraint satisfaction and qualitative reasoning.}
The patchwork property (Theorem~\ref{thm:patchwork}) was established
by Renz and Nebel~\cite{RenzNebel1999} and
Renz~\cite{Renz1999} in the context of qualitative spatial reasoning
and constraint satisfaction.  Our contribution is to bridge the gap
between this CSP result and the description logic setting: we show
that the patchwork property, when combined with the quasimodel method
or a tableau calculus, is precisely the tool needed to decide
satisfiability in $\ALCIRCC{}$ logics.

%% ============================================================
\section{Discussion and Open Questions}
\label{sec:discussion}

\begin{enumerate}
  \item \textbf{Decidability of $\ALCIRCC{5}$ and $\ALCIRCC{8}$.}\;
    Both problems remain \textbf{open}.
    The \PSPACE{} lower bound for $\ALCIRCC{5}$~\cite{Wessel2003}
    and the \EXPTIME{} lower bound for $\ALCIRCC{8}$~\cite{Wessel2003}
    remain valid.  An \EXPTIME{} upper bound was claimed in this paper
    but depends on the type-elimination algorithm
    (Theorem~\ref{thm:complexity}), which is unsound
    (Remark~\ref{rem:type-elim-bug}).  Finding a correct decision
    procedure is the primary open problem.

  \item \textbf{General characterization.}\; Which composition tables
    (beyond RCC5 and RCC8) yield decidable $\ALCIRCC{}$-style logics?
    Our proof shows that the patchwork property is sufficient.  It
    would be valuable to characterize exactly which qualitative
    calculi have this property and whether it is also necessary for
    decidability.

  \item \textbf{Adding number restrictions.}\; Unqualified number
    restrictions make $\ALCRASG$ undecidable~\cite{Wessel2000}.  We
    expect the same for $\ALCIRCC{5}$: number restrictions would allow
    enforcing specific branching patterns that, combined with concept
    distinctions, may enable grid encodings.

  \item \textbf{Finite model reasoning.}\; Satisfiability restricted
    to \emph{finite} models (ruling out infinite PP-chains) may be
    decidable in a lower complexity class, since the need for infinite
    witnesses---and hence the full power of the Henkin
    construction---disappears.

  \item \textbf{Hybrid extensions.}\; Wessel~\cite{Wessel2003}
    considers hybrid logic extensions with nominals and the
    $\mathord{@}$-operator.  Whether our decidability results extend
    to these enrichments is an open question; the interaction between
    nominals and the complete-graph semantics may introduce new
    complications.
\end{enumerate}

%% ============================================================
\section*{Acknowledgments}

This research was prompted by Michael Wessel, who introduced the
$\ALCIRCC{}$ family during his doctoral work at the University of
Hamburg under the DFG project ``Description Logics and Spatial
Reasoning'' (grant NE~279/8-1) and left the decidability of
$\ALCIRCC{5}$ and $\ALCIRCC{8}$ as open problems.  The present
work directly builds on his foundational contributions~\cite{Wessel2003,Wessel2000}.

This revised version corrects several issues identified in a technical
review of the original manuscript: the treatment of $\EQ$ in the
syntax and pair-types (Section~\ref{sec:closure}), the quasimodel
conditions and their extractability from models
(Section~\ref{sec:quasimodel}), the distinction between algebraic
closure and full path-consistency
(Remark~\ref{rem:q3-variants}), the Henkin construction and its
reliance on full RCC5 tractability
(Claim~\ref{claim:extension} and Remark~\ref{rem:extension-gap}),
and the RCC8 extension (Section~\ref{sec:rcc8}).  We are grateful for
the detailed and constructive review.

%% ============================================================
\bibliographystyle{plain}
\begin{thebibliography}{10}

\bibitem{BaaderEtAl2003}
F.~Baader, D.~Calvanese, D.~McGuinness, D.~Nardi, and P.~Patel-Schneider,
  editors.
\newblock {\em The Description Logic Handbook: Theory, Implementation, and
  Applications}.
\newblock Cambridge University Press, 2003.

\bibitem{BorgwardtEtAl2024}
S.~Borgwardt, F.~De~Bortoli, and P.~Koopmann.
\newblock The precise complexity of reasoning in {$\mathcal{ALC}$} with
  $\omega$-admissible concrete domains.
\newblock In {\em Proc.\ KR}, 2024.

\bibitem{BaaderRydval2020}
F.~Baader and M.~Rydval.
\newblock Description logics with concrete domains and general concept
  inclusions revisited.
\newblock In {\em Proc.\ IJCAR}, volume 12166 of LNCS, pages 413--431.
  Springer, 2020.

\bibitem{DemriGu2026}
S.~Demri and T.~Gu.
\newblock Robustness of constraint automata for description logics with
  concrete domains.
\newblock In {\em Proc.\ CSL}, volume 363 of LIPIcs, 2026.

\bibitem{GradelOtto1999}
E.~Gr\"{a}del and M.~Otto.
\newblock On logics with two variables.
\newblock {\em Theoretical Computer Science}, 224:73--113, 1999.

\bibitem{LutzMilicic2007}
C.~Lutz and M.~Mili\v{c}i\'{c}.
\newblock A tableau algorithm for description logics with concrete domains and
  general {TBoxes}.
\newblock {\em Journal of Automated Reasoning}, 38:227--259, 2007.

\bibitem{LutzWolter2006}
C.~Lutz and F.~Wolter.
\newblock Modal logics of topological relations.
\newblock {\em Logical Methods in Computer Science}, 2(2), 2006.

\bibitem{RandellEtAl1992}
D.~Randell, Z.~Cui, and A.~Cohn.
\newblock A spatial logic based on regions and connection.
\newblock In {\em Proc.\ KR}, pages 165--176, 1992.

\bibitem{Renz1999}
J.~Renz.
\newblock Maximal tractable fragments of the {Region Connection Calculus}: A
  complete analysis.
\newblock In {\em Proc.\ IJCAI}, pages 448--454, 1999.

\bibitem{RenzNebel1999}
J.~Renz and B.~Nebel.
\newblock On the complexity of qualitative spatial reasoning: A maximal
  tractable fragment of the {Region Connection Calculus}.
\newblock {\em Artificial Intelligence}, 108(1-2):69--123, 1999.

\bibitem{Wessel2000}
M.~Wessel.
\newblock Decidable and undecidable extensions of {$\mathcal{ALC}$} with
  composition-based role inclusion axioms.
\newblock Technical Report FBI-HH-M-301/01, University of Hamburg, 2000.

\bibitem{Wessel2003}
M.~Wessel.
\newblock Qualitative spatial reasoning with the {$\mathcal{ALCI}_\mathit{RCC}$}
  family---{F}irst results and unanswered questions.
\newblock Technical Report FBI-HH-M-324/03, University of Hamburg, 2002/2003.

\end{thebibliography}

\end{document}
